\documentclass{article}
\usepackage{amsfonts,amssymb,amsmath,amsthm,mathtools,bm}
\numberwithin{equation}{section}
\usepackage{booktabs}
\usepackage{enumitem}
\setlist[itemize]{noitemsep,topsep=2pt}
\usepackage[margin=1in]{geometry}
\usepackage[numbers]{natbib}
\usepackage{authblk}
\usepackage[ruled,linesnumbered]{algorithm2e}
\usepackage{xcolor}
\definecolor{purple}{rgb}{0.50,0.00,0.50}
\usepackage[colorlinks,citecolor=purple,linkcolor=blue]{hyperref}
\usepackage[nameinlink]{cleveref}

\renewcommand{\baselinestretch}{1.15}

\newtheorem{theorem}{Theorem}[section]
\newtheorem{corollary}{Corollary}[section]
\newtheorem{lemma}{Lemma}[section]
\newtheorem{proposition}{Proposition}[section]
\newtheorem{definition}{Definition}[section]
\newtheorem{remark}{Remark}[section]
\newtheorem{assumption}{Assumption}[section]

\crefname{assumption}{Assumption}{Assumptions}
\Crefname{assumption}{Assumption}{Assumptions}
\crefname{algorithm}{Algorithm}{Algorithms}
\Crefname{algorithm}{Algorithm}{Algorithms}

\newcommand{\R}{\mathbb{R}}


\title{Additive Universal Trust-Region Methods\\
for Unconstrained and NC--SC Minimax Optimization}
\author{}
\date{}

\begin{document}
\maketitle

\begin{abstract}
We give three analyses of additive universal trust-region (UTR)
methods.  First, for unconstrained nonconvex optimization, we develop
a certified-inexact adaptive clipped UTR whose outer scale update does
not use a Hessian-Lipschitz constant.  The inexact theorem is stated
under explicit certificate conditions, whereas its exact-oracle
specialization is fully parameter free and has trial scale
$\bar\sigma=\max\{\sigma_0,\gamma_+\sqrt{M/6}\}$; when
$\sigma_0=O(\sqrt M)$, the total number of trust-region subproblems is
\[
 O\!\left(\Delta_F\sqrt M\,\epsilon^{-3/2}
 +\log_+\!\frac{\|\nabla F(x_0)\|}{\epsilon}
 +\log_+\!\frac{\sqrt M}{\sigma_0}+1\right).
\]
Thus all nonconstant search overhead is additive, and the leading term
has the optimal dependence on the initial gap, Hessian smoothness, and
tolerance.  A fixed-scale inexact UTR estimate is developed only as a
technical outer interface and supplies, in addition, the cubic-path
budget needed below.

Second, for nonconvex--strongly-concave minimax optimization, we
separate the outer method from a residual-based inner maximization
module.  Any module returning the prescribed inner residual yields
\[
 O\!\left(\Delta_P\kappa^{3/2}\sqrt\rho\,\epsilon^{-3/2}
 +\log_+\!\frac{\|\nabla P(x_0)\|}{\epsilon}+1\right)
\]
outer iterations and corrected second-order oracle constructions.  If
the module has an accelerated $\sqrt\kappa$ residual-reduction bound,
the leading inner-gradient complexity is
\[
 O\!\left((1+\log\kappa)\Delta_P\kappa^2\sqrt\rho\,
 \epsilon^{-3/2}\right).
\]
A Newton-corrected value-function gradient makes the inner error
quadratic; relative to the explicit raw-gradient analysis of Luo et
al., this removes the leading multiplicative $\log(1/\epsilon)$.

Third, we instantiate the same residual interface with the
parameter-free SCAR solver and give two fully observable outer
constructions.  An over-accurate adaptive SCAR--UTR retains the sharp
outer rate but restores a leading inner $\log(1/\epsilon)$.  An
adaptive fixed-radius variant with persistent SCAR instead has leading
outer and inner complexities
\[
 O(\Delta_P\kappa^{3/2}\sqrt\rho\,\epsilon^{-3/2}),
 \qquad
 O(\Delta_P\kappa^2\sqrt\rho\,\epsilon^{-3/2}),
\]
respectively, with only additive instance-dependent lower-order terms.
Both parameter-free guarantees are asymptotic but class-uniform: after
fixing common regularity constants $(\ell,\mu,\rho)$, one unknown
threshold works for every payoff in the class and every initialization.
For the over-accurate construction, the displayed sufficient threshold
is $\Theta(\mu)$ as $\mu\downarrow0$ with $\ell$ and $\rho$ fixed.  The
algorithms neither use these constants nor certify that the requested
tolerance lies below the relevant threshold.
\end{abstract}

\section{Introduction}

We first consider smooth unconstrained optimization
$\min_x F(x)$.  Approximate first-order stationarity does not exclude
strict saddle points, so the target also controls the smallest Hessian
eigenvalue.  Additive universal trust-region (UTR) models obtain the
standard $\epsilon^{-3/2}$ second-order scale using a quadratic
subproblem whose regularization and radius depend on the gradient
\citep{jiang2026utr}.

Our first contribution is a certified-inexact adaptive UTR that
replaces the unknown scale $\sqrt{M/6}$ in the outer update by a
persistent trial value.
A multiplier-aware acceptance test controls both a certified objective
proxy and gradient growth, while the persistent update amortizes
rejected trials.  Its exact-oracle specialization has sharp leading
dependence on the initial objective gap, Hessian-Lipschitz constant,
and target tolerance, with only additive search overhead.
This is the consistently adaptive behavior emphasized in
\citet{hamad2022consistently} and, for cubic regularization,
\citet{fang2026parameter}.  We also prove a fixed-scale inexact UTR
estimate used only as an outer estimate: uniform
gradient error can make a small interior step increase the objective,
and the proof charges that increase to boundary decrease.

After completing both unconstrained analyses, we consider the
nonconvex--strongly-concave value function
$P(x)=\max_y f(x,y)$.  Luo, Li, and Chen apply an inexact
cubic-regularized Newton method to this problem
\citep{luo2022finding}.  Write $\kappa=\ell/\mu$, let $\rho$ be the
Hessian-Lipschitz constant of $f$, and set
$\Delta_P=P(x_0)-\inf_xP(x)$.  Their outer complexity is
\[
 O\!\left(
  \Delta_P\kappa^{3/2}\sqrt\rho\,\epsilon^{-3/2}
 \right),
\]
and their stated gradient complexity is
$\widetilde O(\Delta_P\kappa^2\sqrt\rho\,\epsilon^{-3/2})$.
In the small-$\epsilon$ branch of their explicit worst-case bound, the
raw estimator $\nabla_x f(x,y)$ enforces an $O(\epsilon/\ell)$ inner
distance and produces a leading $\log(1/\epsilon)$.

The minimax specialization here instead uses a Newton-corrected
value-function gradient.  Its inner error is quadratic, so an
accelerated inner module only needs a distance of order
$\sqrt\epsilon$.  Our three contributions can be summarized as
follows.

\begin{enumerate}[label=(\roman*)]
 \item For unconstrained optimization, the certified-inexact adaptive
 UTR admits value, gradient, and Hessian certificates and can therefore
 be instantiated for the minimax value function whenever such
 certificates are supplied.  As an exact-oracle by-product, the method
 does not use $M$ and, whenever
 $\sigma_0=O(\sqrt M)$, solves at most
 \[
 O\!\left(\Delta_F\sqrt M\,\epsilon^{-3/2}
 +\log_+\!\frac{\|\nabla F(x_0)\|}{\epsilon}
 +\log_+\!\frac{\sqrt M}{\sigma_0}+1\right)
 \]
 trial subproblems.  Its returned point is an
 $O(\epsilon,\sqrt{M\epsilon})$-SOSP, and no logarithm multiplies the
 leading term, in the consistently adaptive sense emphasized by
 \citet{hamad2022consistently,fang2026parameter}.

 \item Gradient-norm regularized trust-region models have already
 been considered for NC--SC minimax problems by
 \citet{wang2025gradient}, while \citet{luo2022finding} analyzes a
 cubic-regularized method.  Our contribution is an inexact-oracle
 reduction based on a Newton-corrected envelope gradient and a modular
 inner residual interface.  It gives
 \[
 O\!\left(\Delta_P\kappa^{3/2}\sqrt\rho\,\epsilon^{-3/2}
 +\log_+\!\frac{\|\nabla P(x_0)\|}{\epsilon}+1\right)
 \]
 corrected second-order oracle constructions and, for an accelerated
 inner module, leading inner-gradient complexity
 \[
 O\!\left((1+\log\kappa)\Delta_P\kappa^2\sqrt\rho\,
 \epsilon^{-3/2}\right).
 \]
 At the level of the proved worst-case bounds, this improves the raw
 plug-in-gradient analysis of \citet{luo2022finding} by one leading
 factor $\log(1/\epsilon)$.

 \item We combine the same residual interface with SCAR to remove all
 problem parameters from the executable minimax algorithms.  The
 over-accurate adaptive SCAR--UTR has sharp leading second-order
 complexity but pays a leading inner $\log(1/\epsilon)$.  The
 fixed-radius persistent-SCAR variant removes that logarithm and has
 leading outer and inner bounds
 \[
 O(\Delta_P\kappa^{3/2}\sqrt\rho\,\epsilon^{-3/2}),
 \qquad
 O(\Delta_P\kappa^2\sqrt\rho\,\epsilon^{-3/2}).
 \]
 These parameter-free results are asymptotic but class-uniform: for
 fixed common regularity constants, one small-tolerance threshold works
 over the entire problem class and all initializations.  For the
 over-accurate construction, the displayed sufficient threshold scales
 as $\Theta(\mu)$ for fixed $\ell$ and $\rho$.  The thresholds are not
 observable, so the results are not a posteriori finite-accuracy
 certificates.
\end{enumerate}

\paragraph{Organization.}
\Cref{sec:preliminaries} introduces the unconstrained problem.
\Cref{sec:generic-utr} records the fixed-scale inexact estimates used
later in the minimax proof, and
\Cref{sec:adaptive-utr} gives a certified-inexact adaptive method and
its complete trial complexity.  \Cref{sec:ncsc} then introduces the
minimax setting, the residual-based inner module, and the corrected
oracle reduction.  \Cref{sec:pf-scar-utr,sec:pf-fixed-radius} give the
two parameter-free SCAR implementations.
\Cref{app:proofs} collects the deferred proofs, and
\Cref{app:parameter-free-scope} explains the class-uniform quantifiers
and the limitation of residual-only certificates.

\section{Unconstrained problem and preliminaries}
\label{sec:preliminaries}

We first consider
\begin{equation}\label{eq:generic-problem}
 \min_{x\in\R^n}F(x),
\end{equation}
where $F^*:=\inf_xF(x)>-\infty$ and
$\Delta_F:=F(x_0)-F^*$.

\begin{assumption}[Hessian regularity]\label{ass:generic}
The function $F$ is twice continuously differentiable and, for some
$M>0$,
\begin{equation}\label{eq:hessian-lipschitz}
 \|\nabla^2F(x)-\nabla^2F(z)\|
 \le M\|x-z\|
 \qquad\forall x,z\in\R^n.
\end{equation}
\end{assumption}

\subsection{Notation and approximate stationarity}

The norm $\|\cdot\|$ is Euclidean for vectors and spectral for matrices.
For a symmetric matrix $H$, $\lambda_{\min}(H)$ denotes its smallest
eigenvalue.  We write
$\log_+(t)=\max\{\log t,0\}$ for $t>0$ and set $\log_+(0)=0$.
We also write $\mathbb N_0:=\{0,1,2,\ldots\}$.
The notation $\widetilde O(\cdot)$ suppresses logarithmic factors.

\begin{definition}[Approximate stationary points]\label{def:stationarity}
For a twice differentiable function $F$ and tolerances
$\epsilon,\epsilon_g,\epsilon_H>0$, a point $x$ is an
$\epsilon$-first-order stationary point if
$\|\nabla F(x)\|\le\epsilon$.  It is an
$(\epsilon_g,\epsilon_H)$-second-order stationary point if
\[
 \|\nabla F(x)\|\le\epsilon_g,
 \qquad
 \lambda_{\min}(\nabla^2F(x))\ge-\epsilon_H.
\]
Our generic target is an
$(\epsilon,\sqrt{M\epsilon})$-second-order stationary point when
$\nabla^2F$ is $M$-Lipschitz.
\end{definition}

We repeatedly use the following standard consequences of Hessian
Lipschitz continuity.

\begin{lemma}[Second-order Taylor bounds]\label{lem:taylor-bounds}
If $\nabla^2F$ is $M$-Lipschitz, then, for all $x,d$,
\begin{subequations}\label{eq:taylor-bounds}
\begin{equation}\label{eq:gradient-taylor}
 \bigl\|\nabla F(x+d)-\nabla F(x)-\nabla^2F(x)d\bigr\|
 \le \frac M2\|d\|^2.
\end{equation}
\begin{equation}\label{eq:function-taylor}
 F(x+d)\le{}F(x)+\nabla F(x)^Td
 +\frac12d^T\nabla^2F(x)d+\frac M6\|d\|^3.
\end{equation}
\end{subequations}
These estimates are standard; see
\citet[Lemma~1.2.4]{nesterov2004introductory}.
\end{lemma}

\begin{lemma}[Global optimality for a quadratic trust-region problem]
\label{lem:tr-global-optimality}
Let $g\in\R^n$, let $B\in\R^{n\times n}$ be symmetric, and let
$R>0$.  A vector $d$ is a global solution of
\[
 \min_{\|u\|\le R}\;g^Tu+\frac12u^TBu
\]
if and only if there is a multiplier $\lambda\ge0$ such that
\begin{equation}\label{eq:tr-global-optimality}
 (B+\lambda I)d=-g,\qquad
 B+\lambda I\succeq0,\qquad
 \|d\|\le R,\qquad
 \lambda(\|d\|-R)=0.
\end{equation}
In particular, $\lambda>0$ implies $\|d\|=R$.
\end{lemma}

This characterization is classical; see, for example,
\citet[Section~8.7]{luenberger2021linear}.  All trust-region
subproblems below are solved globally, so
\Cref{lem:tr-global-optimality} applies throughout.


\section{Fixed-scale inexact UTR estimates}
\label{sec:generic-utr}
This section isolates the fixed-scale estimates used later in the
minimax reduction.  These estimates analyze repeated, globally solved
additive UTR subproblems; they are not presented as a separate
algorithmic contribution.  All technical proofs are deferred to
\Cref{app:fixed-scale-proofs}.

\begin{assumption}[Uniform derivative accuracy]
\label{ass:inexact-oracle}
Fix a target tolerance $\epsilon>0$.  At every iterate, the inexact
oracle returns
$g_k\in\R^n$ and a symmetric $H_k\in\R^{n\times n}$ satisfying
\begin{equation}\label{eq:generic-oracle-errors}
  \|g_k-\nabla F(x_k)\|\le\frac{\epsilon}{256},
  \qquad
  \|H_k-\nabla^2F(x_k)\|
  \le\frac1{64}\sqrt{\frac M6}\sqrt\epsilon.
\end{equation}
\end{assumption}
Set
\begin{equation}\label{eq:generic-scale}
  h_k:=\max\{\|g_k\|,\epsilon\}.
\end{equation}
At iterate $x_k$, terminate if
\begin{equation}\label{eq:generic-stopping-test}
 \|g_k\|\le\epsilon,
 \qquad
 \lambda_{\min}(H_k)\ge-\sqrt{\frac M6}\sqrt\epsilon.
\end{equation}
Otherwise, compute a global solution $d_k$ of
\begin{equation}\label{eq:generic-subproblem}
\begin{aligned}
  \min_{d\in\R^n}\quad
  &g_k^Td+\frac12d^T\left(H_k+\sqrt{\frac M6}\sqrt{h_k}I\right)d,\\
  \operatorname{s.t.}\quad
  &\|d\|\le\frac{\sqrt{h_k}}{4\sqrt{M/6}},
\end{aligned}
\end{equation}
and set $x_{k+1}=x_k+d_k$.  This is the additive UTR subproblem of
\citet{jiang2026utr}, evaluated with the derivative estimates
$(g_k,H_k)$; the trust-region subproblem itself is solved exactly.

\subsection{Complexity estimates}

The one-step descent, clipped-gradient recurrence, and finite-prefix
counting estimates are stated and proved in
\Cref{app:fixed-scale-proofs}.  Their two consequences needed in the
minimax analysis are recorded next.

\begin{theorem}[Fixed-scale inexact UTR estimate]
\label{thm:generic-complexity}
Fix $\epsilon>0$.  Suppose $F^*>-\infty$ and
\Cref{ass:generic,ass:inexact-oracle} hold.  Starting from $x_0$,
repeatedly apply the stopping rule
\eqref{eq:generic-stopping-test} and the globally solved subproblem
\eqref{eq:generic-subproblem}.  This procedure terminates after
\begin{equation}\label{eq:generic-complexity}
  O\!\left(
    \Delta_F\sqrt M\,\epsilon^{-3/2}
    +\log_+\!\frac{h_0}{\epsilon}+1
  \right)
\end{equation}
iterations and trust-region subproblem solves.  Its output $\widehat x$
satisfies
\begin{equation}\label{eq:generic-return}
  \|\nabla F(\widehat x)\|\le\frac{257}{256}\epsilon,
  \qquad
  \lambda_{\min}(\nabla^2F(\widehat x))
  \ge-\sqrt{M\epsilon}.
\end{equation}
Replacing the internal tolerance by $256\epsilon/257$ gives the
standard $(\epsilon,\sqrt{M\epsilon})$ guarantee without changing any
complexity order.
\end{theorem}




\begin{lemma}[Cubic path budget]\label{lem:cubic-path}
Fix $\epsilon>0$ and suppose that \Cref{ass:generic,ass:inexact-oracle}
hold.  Let $K$ and $r_k=\|d_k\|$ be generated by
\eqref{eq:generic-stopping-test}--\eqref{eq:generic-subproblem}, where
$K<\infty$ is the return index guaranteed by
\Cref{thm:generic-complexity}.  Then
\begin{equation}\label{eq:cubic-path}
  M\sum_{k=0}^{K-1}r_k^3
  \le\frac{384}{31}\Delta_F
    +\frac{K}{31\sqrt{62}}\frac{\epsilon^{3/2}}{\sqrt M}.
\end{equation}
\end{lemma}




\section{Certified-inexact adaptive clipped UTR}
\label{sec:adaptive-utr}
The following adaptive method does not require the value of $M$ as an
algorithmic input.  It accepts certified approximations of the
objective, gradient, and Hessian; exact oracles are the zero-error
special case.  This formulation is essential for the NC--SC
specialization in \Cref{sec:ncsc}, where the value function is not
evaluated exactly.  The oracle cost required to produce a requested
certificate is kept separate from the number of certified trials.

\begin{remark}[Why the fixed-scale estimate remains separate]
\label{rem:fixed-not-adaptive-special-case}
The fixed-scale result in \Cref{sec:generic-utr} and the adaptive result
below share the trust-region geometry, but the former is not a formal
special case of the latter.  The fixed-scale procedure accepts every
step without querying function values, permits a controlled increase
on an inexact interior step, and supplies the cubic-path budget
\eqref{eq:cubic-path}.  The adaptive method instead filters every step
through certified value and gradient tests and uses a padded Hessian
model.  Consequently, specializing the adaptive theorem would neither
justify the unconditional fixed-scale update nor recover its
cubic-path estimate.  We keep the two counting arguments separate,
but derive both one-step bounds from the common estimate in
\Cref{lem:common-inexact-tr-step}; that estimate, in turn, uses the
single trust-region characterization in
\Cref{lem:tr-global-optimality}.
\end{remark}

\begin{assumption}[Certified adaptive oracle and parameter regime]
\label{ass:adaptive-inexact-oracle}
Suppose \Cref{ass:generic} holds and $F^*>-\infty$.  Fix
\[
 \epsilon>0,\qquad
 0<\sigma_{\min}\le\sigma_0,\qquad
 \gamma_+>1,\qquad \gamma_->1,\qquad
 0<\eta\le\frac1{256},
\]
together with a value radius $\omega\ge0$ and a gradient radius
$0\le\delta_g\le\epsilon/256$.  Define
\begin{equation}\label{eq:adaptive-scale-bar}
 \sigma_*:=\sqrt{\frac M6},
 \qquad
 \bar\sigma:=\max\{\sigma_0,\gamma_+\sigma_*\},
\end{equation}
and assume
\begin{subequations}\label{eq:adaptive-safe-compatibility}
\begin{align}
 \frac{\delta_g^2}{\sigma_*\sqrt\epsilon}+\frac\omega4
 &\le\omega,
 \label{eq:adaptive-interior-compatibility}\\
 \left(\frac{15}{1024}-\eta\right)
 \frac{\epsilon^{3/2}}{\bar\sigma}
 &\ge\frac{5\omega}{4}.
 \label{eq:adaptive-boundary-compatibility}
\end{align}
\end{subequations}
The associated dimensionless inexactness load satisfies
\begin{equation}\label{eq:adaptive-absorption-load}
 \chi:=\left(1+\frac4{\log2}\right)
       \frac{\bar\sigma\omega}{\eta\epsilon^{3/2}}<1.
\end{equation}

At a current iterate $x_k$, a value--gradient query returns a stored
pair $(\widetilde F_k,g_k)$ satisfying
\begin{equation}\label{eq:adaptive-vg-errors}
 |\widetilde F_k-F(x_k)|\le\frac\omega8,
 \qquad
 \|g_k-\nabla F(x_k)\|\le\delta_g.
\end{equation}
At a candidate $x_k^+=x_k+d$, the notation
$(\widetilde F_k^+,g_k^+)$ denotes a pair satisfying the same bounds
with $x_k$ replaced by $x_k^+$.  For every trial scale $\sigma>0$, a
Hessian query at $x_k$ returns a symmetric $H_{k,\sigma}$ satisfying
\begin{equation}\label{eq:adaptive-hessian-error}
 \|H_{k,\sigma}-\nabla^2F(x_k)\|
 \le\frac{\sigma\sqrt\epsilon}{64}.
\end{equation}
An accepted candidate value--gradient pair is stored without
recomputation.  A new Hessian may be requested at the next trial
scale.
\end{assumption}

At an outer iterate $x_k$, write
\begin{equation}\label{eq:clipped-scale}
 h_k:=\max\{\|g_k\|,\epsilon\}.
\end{equation}
For a trial value $\sigma>0$, consider
\begin{equation}\label{eq:adaptive-subproblem}
\begin{aligned}
    \min_{d\in\R^n}\quad
        &g_k^Td+\frac12d^T\left(
          H_{k,\sigma}+\frac{\sigma\sqrt\epsilon}{64}I
          +\sigma\sqrt{h_k}I\right)d,\\
    \operatorname{s.t.}\quad
        &\|d\|\le\frac{\sqrt{h_k}}{4\sigma}.
\end{aligned}
\end{equation}
The subproblem is solved globally.  Its trust-region multiplier is denoted
by $\lambda\ge0$, as supplied by
\Cref{lem:tr-global-optimality} with
$B=H_{k,\sigma}+(\sigma\sqrt\epsilon/64+\sigma\sqrt{h_k})I$ and
$R=\sqrt{h_k}/(4\sigma)$.  In particular, $\lambda>0$ implies that the
radius is active.

The trial is accepted if one of the following mutually exclusive tests
holds:
\begin{subequations}\label{eq:acceptance}
\begin{align}
    \|d\|<\frac{\sqrt{h_k}}{4\sigma}:\quad
    &\widetilde F_k^+\le \widetilde F_k+\omega,
    \qquad
    \|g_k^+\|\le\frac12h_k;
       \label{eq:acceptance-interior}\\
    \|d\|=\frac{\sqrt{h_k}}{4\sigma}:\quad
    &\widetilde F_k-\widetilde F_k^+
       \ge\eta\left(
           \frac{h_k^{3/2}}{\sigma}+\lambda\|d\|^2
       \right)+\omega, \label{eq:acceptance-boundary-function}\\
    &\|g_k^+\|\le\frac12h_k+\lambda\|d\|.
       \label{eq:acceptance-boundary-gradient}
\end{align}
\end{subequations}

\begin{algorithm}[H]
\DontPrintSemicolon
\caption{Certified-inexact adaptive clipped UTR}
\label{alg:adaptive-utr}
\KwIn{$x_0$ and parameters and certified oracle responses satisfying
      \Cref{ass:adaptive-inexact-oracle}}
$k\leftarrow0$, $\sigma\leftarrow\sigma_0$\;
obtain and store $(\widetilde F_k,g_k)$ at $x_0$\;
\While{true}{
    $h_k\leftarrow\max\{\|g_k\|,\epsilon\}$\;
    \Repeat{a step is accepted}{
        obtain $H_{k,\sigma}$ at $(x_k,\sigma)$ satisfying
        \eqref{eq:adaptive-hessian-error}\;
        \If{$\|g_k\|\le\epsilon$ and
             $\lambda_{\min}(H_{k,\sigma})
              \ge-(9/8)\sigma\sqrt\epsilon$}{
            \Return{$x_k$}\;
        }
        Solve \eqref{eq:adaptive-subproblem} and obtain $(d,\lambda)$\;
        obtain $(\widetilde F_k^+,g_k^+)$ at $x_k+d$\;
        \eIf{\eqref{eq:acceptance} holds}{
            $x_{k+1}\leftarrow x_k+d$\;
            $(\widetilde F_{k+1},g_{k+1})
             \leftarrow(\widetilde F_k^+,g_k^+)$\;
            $\sigma\leftarrow\max\{\sigma_{\min},\sigma/\gamma_-\}$\;
            $k\leftarrow k+1$\;
            accept the step\;
        }{
            discard $(\widetilde F_k^+,g_k^+)$\;
            $\sigma\leftarrow\gamma_+\sigma$\;
        }
    }
}
\end{algorithm}

The stopping test in Algorithm~\ref{alg:adaptive-utr} is evaluated for
every trial value of $\sigma$, including immediately after each rejected
trial increases $\sigma$.

\subsection{Trial-step estimates}
The safe-scale, parameter-bound, and finite-prefix counting lemmas used
to prove the main theorem are stated and proved in
\Cref{app:proofs}.
The supporting estimates in \Cref{app:proofs} establish finite
termination and the following main complexity guarantee.

\begin{remark}[Parameter freedom versus certified inexactness]
\label{rem:adaptive-certificate-parameters}
The compatibility conditions in
\Cref{ass:adaptive-inexact-oracle} are analytical sufficient
conditions, not tests performed by the algorithm.  For a nonzero
certificate radius, their verification may use $M$ through
$\sigma_*$ and $\bar\sigma$.  Thus the certified-inexact theorem is
conditional on an oracle provider that can supply the stated radii; it
does not claim that producing such certificates is parameter free.
This does not affect the first contribution: with exact oracle
responses, $\omega=\delta_g=0$, all compatibility conditions hold
automatically and the executable scale update never uses $M$.
\end{remark}

\begin{theorem}[Certified-inexact adaptive UTR complexity]
\label{thm:total-complexity}
Suppose \Cref{ass:adaptive-inexact-oracle} holds and define
$\Delta_F:=F(x_0)-F^*$.  Then
\Cref{alg:adaptive-utr} terminates after finitely many trials.  If
$Q_\epsilon$ denotes the total number of solved trial subproblems, then
\begin{equation}\label{eq:adaptive-total-big-o}
 Q_\epsilon=O\!\left(
 \frac{\bar\sigma(\Delta_F+\omega)\epsilon^{-3/2}
       +\log_+(h_0/\epsilon)+1}{1-\chi}
 +\log_+\!\frac{\bar\sigma}{\sigma_0}
 \right),
\end{equation}
where the hidden constant depends only on $\eta$ and
$\gamma_\pm$.  The returned point $\widehat x$ satisfies
\begin{equation}\label{eq:return-guarantee}
    \|\nabla F(\widehat x)\|\le\epsilon+\delta_g,
    \qquad
    \lambda_{\min}(\nabla^2 F(\widehat x))
    \ge-\frac{73}{64}\bar\sigma\sqrt\epsilon.
\end{equation}
\end{theorem}




\begin{remark}[Exact-oracle specialization and the initial scale]
\label{rem:exact-oracle-overhead}
Suppose $F^*>-\infty$ and \Cref{ass:generic} holds, and run
\Cref{alg:adaptive-utr} with exact function values, gradients, and
Hessians.  Then $\omega=\delta_g=0$, so the compatibility conditions in
\Cref{ass:adaptive-inexact-oracle} hold automatically.  No restriction
on $\sigma_0$ is needed for the general bound
\begin{equation}\label{eq:big-O-complexity}
 Q_\epsilon=O\!\left(
  \bar\sigma\Delta_F\epsilon^{-3/2}
  +\log_+(h_0/\epsilon)
  +\log_+(\bar\sigma/\sigma_0)+1
 \right),
 \qquad
 \bar\sigma=\max\{\sigma_0,\gamma_+\sqrt{M/6}\},
\end{equation}
and the return guarantee is
$\lambda_{\min}(\nabla^2F(\widehat x))
\ge-(73/64)\bar\sigma\sqrt\epsilon$.
The additional condition $\sigma_0=O(\sqrt M)$ is needed only when
one wants to simplify $\bar\sigma=O(\sqrt M)$ and obtain the sharp
leading term $O(\sqrt M\,\Delta_F\epsilon^{-3/2})$ together with an
$O(\sqrt{M\epsilon})$ second-order guarantee.  Without this
initial-scale condition, the general bound and the
$O(\bar\sigma\sqrt\epsilon)$ curvature guarantee above remain valid.
\end{remark}

The theorem counts certified trials and globally solved trust-region
subproblems.  In the exact-oracle specialization, the Hessian is evaluated
once per trial and the candidate function and gradient once per solved
subproblem.  For an inexact implementation, the work needed to produce
the certificates in \Cref{ass:adaptive-inexact-oracle} must be counted
separately.
The persistent
parameter update in \Cref{alg:adaptive-utr} is the same amortization
mechanism used in adaptive regularization methods
\citep{cartis2011adaptive2}; the multiplier-aware acceptance test is
what couples gradient growth to objective decrease.

\section{NC--SC minimax optimization}\label{sec:ncsc}

\subsection{Problem setting and value-function geometry}

We first use the fixed-scale estimates of \Cref{sec:generic-utr} for
\begin{equation}\label{eq:ncsc-problem}
 \min_{x\in\R^{d_x}}P(x),
 \qquad
 P(x):=\max_{y\in\R^{d_y}}f(x,y).
\end{equation}
The certified-inexact adaptive result in
\Cref{thm:total-complexity} is also directly applicable once the
corrected value and derivative certificates below are supplied.  We
retain the fixed-scale wrapper for the explicit total inner-work bound,
because its cubic-path estimate is the mechanism that removes the
leading $\log(1/\epsilon)$.

The minimax development has three layers.  This section first derives
the value-function geometry and a corrected oracle, and then analyzes a
known-parameter outer method whose maximization routine is supplied as
an abstract residual module.  \Cref{sec:pf-scar-utr} removes the
parameters by using uniformly over-accurate SCAR calls, while
\Cref{sec:pf-fixed-radius} combines a fixed-radius outer method with
persistent SCAR to remove the resulting leading logarithms from the
inner-oracle bound.  Keeping the residual interface common to all three
layers makes clear which conclusions use only correctness of the inner
solve and which additionally use a particular work guarantee.

We impose the assumptions used by \citet{luo2022finding}.

\begin{assumption}[NC--SC regularity]\label{ass:ncsc}
The payoff $f$ is twice continuously differentiable and satisfies:
\begin{enumerate}[label=(\roman*)]
  \item $\nabla f$ is $\ell$-Lipschitz on
  $\R^{d_x}\times\R^{d_y}$;
  \item $\nabla^2f$ is $\rho$-Lipschitz on
  $\R^{d_x}\times\R^{d_y}$ for some $\rho>0$;
  \item $f(x,\cdot)$ is $\mu$-strongly concave for every $x$, where
  $0<\mu\le\ell$;
  \item $P^*:=\inf_xP(x)>-\infty$.
\end{enumerate}
We write $\kappa:=\ell/\mu$.
\end{assumption}

Strong concavity makes
$y^*(x):=\arg\max_y f(x,y)$ unique and guarantees that
$\nabla^2_{yy}f(x,y)$ is nonsingular.  Define the Schur complement
\begin{equation}\label{eq:schur-definition}
 \mathcal H(x,y):=\nabla^2_{xx}f(x,y)
 -\nabla^2_{xy}f(x,y)
  \bigl(\nabla^2_{yy}f(x,y)\bigr)^{-1}
  \nabla^2_{yx}f(x,y).
\end{equation}

The main departure from \citet{luo2022finding} is the value-function
gradient estimator.  Their MCN method uses the raw plug-in gradient
$\nabla_xf(x,y)$, whose error is linear in $\|y-y^*(x)\|$, together
with the same Schur-complement Hessian \eqref{eq:schur-definition}.
We instead use the Newton-corrected gradient introduced below, whose
inner error is quadratic.  This relaxation permits an inner accuracy
of order $\sqrt\epsilon$ and an observable residual stopping rule;
at the level of the proved upper bounds, it preserves the leading outer
rate and removes the multiplicative $\log(1/\epsilon)$ from the leading
inner-gradient complexity.

\begin{lemma}[Value-function geometry]\label{lem:value-geometry}
Under \Cref{ass:ncsc}, $y^*(\cdot)$ is $\kappa$-Lipschitz, $P$ is twice
continuously differentiable, and
\begin{subequations}\label{eq:value-derivatives}
\begin{align}
 \nabla P(x)&=\nabla_xf(x,y^*(x)),\label{eq:value-gradient}\\
 \nabla^2P(x)&=\mathcal H(x,y^*(x)).\label{eq:value-hessian}
\end{align}
\end{subequations}
Moreover, $\nabla P$ is $(1+\kappa)\ell$-Lipschitz and
\begin{equation}\label{eq:value-hessian-lipschitz}
 \|\nabla^2P(x)-\nabla^2P(x')\|
 \le M_P\|x-x'\|,
 \qquad
 M_P:=4\sqrt2\,\kappa^3\rho.
\end{equation}
These properties follow from
\citet[Lemmas~1--3]{luo2022finding}.
\end{lemma}

In view of \eqref{eq:value-hessian-lipschitz}, the target in this
section is an
$(\epsilon,\sqrt{M_P\epsilon})$-second-order stationary point of $P$.

\subsection{A residual-based inner maximization module}

The outer method uses the inner maximization only through an observable
residual interface.  This separates outer correctness from the choice
between AGD, SCAR, or another accelerated inner solver.

\begin{lemma}[Residual--distance certificate]
\label{lem:inner-residual-certificate}
Under \Cref{ass:ncsc}, for every $x$ and $y$,
\begin{equation}\label{eq:inner-residual-certificate}
 \|y-y^*(x)\|
 \le\frac{\|\nabla_yf(x,y)\|}{\mu}
 \le\kappa\|y-y^*(x)\|.
\end{equation}
\end{lemma}




\begin{definition}[Inner residual module]
\label{def:inner-residual-module}
Given an outer point $x$, a warm start $y^{\rm in}$, and a target
$\tau>0$, an inner residual module $\mathcal S$ returns
\begin{equation}\label{eq:maxres-interface}
 y^{\rm out}
 =\operatorname{MaxRes}_{\mathcal S}(x,y^{\rm in},\tau)
 \quad\text{such that}\quad
 \|\nabla_yf(x,y^{\rm out})\|\le\tau.
\end{equation}
The module may retain implementation-specific state between calls, but
this state is hidden from the outer method.  We call the module
\emph{accelerated} if there are implementation factors
$C_{\mathcal S}>0$ and $b_{\mathcal S}\ge1$, which may depend on the
implementation and on $\kappa$ but not on $x$, $\tau$, $R_{\rm in}$,
$\epsilon$, or the call index, such that a call whose
initial residual is
$R_{\rm in}:=\|\nabla_yf(x,y^{\rm in})\|$ uses at most
\begin{equation}\label{eq:inner-module-work}
 C_{\mathcal S}\sqrt\kappa\left[
 b_{\mathcal S}
 +\left\lceil\log_2\max\left\{1,
       \frac{R_{\rm in}}{\tau}\right\}\right\rceil
 \right]
\end{equation}
evaluations of $\nabla_yf$.  The terminal residual evaluation may be
reused when constructing the corrected gradient.
\end{definition}

The interface \eqref{eq:maxres-interface}, rather than the work bound
\eqref{eq:inner-module-work}, is the only inner-solver property used by
the outer convergence theorem.  We next record AGD as one
implementation.  For a fixed $x$, starting from $y^{(0)}$, set
$v^{(0)}=y^{(0)}$ and
\begin{equation}\label{eq:agd-recursion}
\begin{aligned}
 y^{(j+1)}&=v^{(j)}+\frac1\ell\nabla_yf(x,v^{(j)}),\\
 v^{(j+1)}&=y^{(j+1)}
 +\frac{\sqrt\kappa-1}{\sqrt\kappa+1}
   \bigl(y^{(j+1)}-y^{(j)}\bigr).
\end{aligned}
\end{equation}

\begin{proposition}[AGD implementation]
\label{prop:agd-implementation}
Under \Cref{ass:ncsc}, fix $x$, $y^{(0)}$, and $\tau>0$.  For
$\kappa>1$, the iterate in \eqref{eq:agd-recursion} satisfies
\begin{equation}\label{eq:agd-contraction}
 \|y^{(N)}-y^*(x)\|
 \le\sqrt{\kappa+1}
 \left(1-\frac1{\sqrt\kappa}\right)^{N/2}
 \|y^{(0)}-y^*(x)\|.
\end{equation}
Consequently, after evaluating
$R_{\rm in}=\|\nabla_yf(x,y^{(0)})\|$, return immediately if
$R_{\rm in}\le\tau$; otherwise run
\[
 N:=\left\lceil2\sqrt\kappa\log_+\!\left(
  \frac{\sqrt{\kappa+1}\,\kappa R_{\rm in}}{\tau}
 \right)\right\rceil
\]
AGD steps and verify the residual once at the terminal iterate.  The
initial gradient is reused in the first AGD step.  This implements
\Cref{def:inner-residual-module} using at most
\begin{equation}\label{eq:agd-module-work}
 1+\left\lceil2\sqrt\kappa\log_+\!\left(
  \frac{\sqrt{\kappa+1}\,\kappa R_{\rm in}}{\tau}
 \right)\right\rceil
\end{equation}
inner-gradient evaluations.  Hence it satisfies
\eqref{eq:inner-module-work} with universal $C_{\mathcal S}$ and
$b_{\mathcal S}=O(1+\log\kappa)$.  When $\kappa=1$, one gradient step
is sufficient.
\end{proposition}




We count one evaluation of $\nabla_yf(x,y)$ as one inner-gradient
oracle call.  Corrected-gradient and Schur-complement constructions,
trust-region subproblem solves, and any linear algebra used to apply
$(\nabla^2_{yy}f)^{-1}$ are reported separately.

\subsection{Corrected oracle and minimax algorithm}

We now construct certified value and derivative oracles for the value
function from derivatives of $f$.  At an approximate maximizer $y$,
define the Newton-corrected value
\begin{equation}\label{eq:corrected-value-pf}
 \widehat P(x,y):=f(x,y)
 -\frac12\nabla_y f(x,y)^T
  \bigl(\nabla^2_{yy}f(x,y)\bigr)^{-1}
  \nabla_y f(x,y).
\end{equation}
Since $\nabla^2_{yy}f\prec0$, the correction raises the raw payoff
and is exact for a concave quadratic inner problem.  The corrected
gradient and Hessian estimators are
\begin{subequations}\label{eq:corrected-derivatives}
\begin{align}
  \widehat g(x,y)
  &:=\nabla_xf(x,y)-\nabla^2_{xy}f(x,y)
       \bigl(\nabla^2_{yy}f(x,y)\bigr)^{-1}\nabla_yf(x,y),
       \label{eq:corrected-gradient}\\
  \widehat H(x,y)
  &:=\mathcal H(x,y).
       \label{eq:corrected-hessian}
\end{align}
\end{subequations}
The gradient estimator is the oracle-level distinction from
\citet[Algorithm~2]{luo2022finding}: their method uses
$\nabla_xf(x,y)$, whereas \eqref{eq:corrected-gradient} adds a Newton
correction.  Both methods use the Schur-complement Hessian
$\mathcal H(x,y)$.

\begin{lemma}[Corrected value and derivative errors]
\label{lem:corrected-errors}
Under \Cref{ass:ncsc}, for every $x$ and $y$,
\begin{subequations}\label{eq:corrected-error-bounds}
\begin{align}
 |\widehat P(x,y)-P(x)|
 &\le\frac{(3\kappa-1)\rho}{12}\|y-y^*(x)\|^3
 \le\frac{\kappa\rho}{4}\|y-y^*(x)\|^3,
       \label{eq:corrected-value-error-pf}\\
 \|\widehat g(x,y)-\nabla P(x)\|
 &\le\frac{(1+\kappa)\rho}{2}\|y-y^*(x)\|^2,
       \label{eq:corrected-gradient-error}\\
 \|\widehat H(x,y)-\nabla^2P(x)\|
 &\le(1+\kappa)^2\rho\|y-y^*(x)\|.
       \label{eq:corrected-hessian-error}
\end{align}
\end{subequations}
Moreover, whenever
$\|y-y^*(x)\|\le\mu/\rho$, the Newton-corrected value satisfies the
following $\kappa$-independent local bound:
\begin{equation}\label{eq:corrected-value-local-error}
 |\widehat P(x,y)-P(x)|
 \le\frac{7\rho}{24}\|y-y^*(x)\|^3.
\end{equation}
\end{lemma}




Set the outer residual target
\begin{equation}\label{eq:outer-residual-target}
 \tau_{\rm out}:=
 \frac{\mu\sqrt{M_P\epsilon/6}}
      {64(1+\kappa)^2\rho}.
\end{equation}
A further practical distinction from
\citet[Algorithm~2 and Theorem~2]{luo2022finding} is the inner
termination rule.  Their analysis prescribes an iteration count for
each AGD solve, whereas our outer method asks the input module
$\mathcal S$ to satisfy the directly evaluable condition
$\|\nabla_yf(x,y)\|\le\tau_{\rm out}$.  By
\eqref{eq:inner-residual-certificate}, this guarantees
$\|y-y^*(x)\|\le\tau_{\rm out}/\mu$.  The resulting oracle
certificates are recorded explicitly next.

\begin{corollary}[Residual-to-oracle certificate]
\label{cor:ncsc-fixed-oracle}
Under \Cref{ass:ncsc}, let
$M_P=4\sqrt2\,\kappa^3\rho$ and let $\tau_{\rm out}$ be given by
\eqref{eq:outer-residual-target}.  Every point $y$ satisfying
$\|\nabla_yf(x,y)\|\le\tau_{\rm out}$ also satisfies
\begin{equation}\label{eq:ncsc-fixed-oracle-errors}
 \|\widehat g(x,y)-\nabla P(x)\|<\frac{\epsilon}{8192},
 \qquad
 \|\widehat H(x,y)-\nabla^2P(x)\|
 \le\frac1{64}\sqrt{\frac{M_P}{6}}\sqrt\epsilon.
\end{equation}
In particular, the fixed-scale inexact oracle conditions
\eqref{eq:generic-oracle-errors} hold with $F=P$ and $M=M_P$.
\end{corollary}

Thus the stopping decision requires neither $y^*(x)$ nor an a priori
bound on its distance from the inner initializer.
The target has the useful scale and lower bound
\begin{subequations}\label{eq:outer-residual-scale}
\begin{align}
 \frac{\tau_{\rm out}}\mu
 &=\Theta\!\left(\sqrt{\frac{\epsilon}{\kappa\rho}}\right)
 =\Theta\!\left(\kappa\sqrt{\frac{\epsilon}{M_P}}\right),
 \label{eq:outer-residual-distance-scale}\\
 \tau_{\rm out}
 &\ge\frac{\ell}{64\sqrt3}
       \sqrt{\frac\epsilon{M_P}}.
 \label{eq:outer-residual-lower}
\end{align}
\end{subequations}
The second inequality follows by substitution: the ratio of its
left-hand side to its right-hand side is
$4\kappa^2/(1+\kappa)^2\ge1$.

\begin{algorithm}[H]
\DontPrintSemicolon
\caption{Fixed-scale NC--SC UTR with an inner residual module}
\label{alg:ncsc-utr}
\KwIn{$x_0,y_{-1}$, $\epsilon>0$, $\ell,\mu,\rho$, and an inner
      module $\mathcal S$ satisfying \eqref{eq:maxres-interface}}
$\kappa\leftarrow\ell/\mu$,
$M_P\leftarrow4\sqrt2\,\kappa^3\rho$,
and compute $\tau_{\rm out}$ from
\eqref{eq:outer-residual-target}\;
$y_0\leftarrow
 \operatorname{MaxRes}_{\mathcal S}(x_0,y_{-1},\tau_{\rm out})$\;
\For{$k=0,1,\ldots$}{
  Form $g_k=\widehat g(x_k,y_k)$ and
  $H_k=\widehat H(x_k,y_k)$\;
  $h_k\leftarrow\max\{\|g_k\|,\epsilon\}$\;
  \If{$\|g_k\|\le\epsilon$ and
      $\lambda_{\min}(H_k)\ge-\sqrt{M_P/6}\sqrt\epsilon$}{
    \Return{$x_k$}\;
  }
  Solve \eqref{eq:generic-subproblem} globally with $F=P$ and
  $M=M_P$, and obtain $d_k$\;
  $x_{k+1}\leftarrow x_k+d_k$\;
  $y_{k+1}\leftarrow
   \operatorname{MaxRes}_{\mathcal S}
   (x_{k+1},y_k,\tau_{\rm out})$\;
}
\end{algorithm}

Thus the minimax algorithm never evaluates $P(x)$ and never solves the
inner maximization to $O(\epsilon)$ distance.  It uses only the residual
contract \eqref{eq:maxres-interface}, with a target whose associated
distance is proportional to $\sqrt\epsilon$.  The outer algorithm and
its correctness proof are independent of how $\mathcal S$ attains that
residual.  The parameters in \eqref{eq:outer-residual-target} and the
outer scale $M_P$ are still known in this baseline; using a
parameter-free inner implementation alone does not make the full
wrapper parameter-free.

\subsection{Outer and inner complexity}

\begin{theorem}[NC--SC complexity under an inner residual module]
\label{thm:ncsc-complexity}
Suppose $f$ satisfies \Cref{ass:ncsc} and define
$P(x)=\max_y f(x,y)$.  Fix
$(x_0,y_{-1})$ and $\epsilon>0$, and set
\[
 \kappa=\ell/\mu,
 \qquad M_P=4\sqrt2\,\kappa^3\rho,
 \qquad \Delta_P=P(x_0)-\inf_xP(x).
\]
Run \Cref{alg:ncsc-utr} with the residual target
\eqref{eq:outer-residual-target} and any inner module that returns a
point with residual at most that target.  Then the algorithm returns
$\widehat x$ satisfying
\begin{equation}\label{eq:ncsc-return}
 \|\nabla P(\widehat x)\|\le\frac{257}{256}\epsilon,
 \qquad
 \lambda_{\min}(\nabla^2P(\widehat x))
 \ge-\sqrt{M_P\epsilon}.
\end{equation}
The numbers of outer iterations, corrected-gradient and
Schur-complement constructions, and trust-region subproblem solves are
bounded by
\begin{equation}\label{eq:ncsc-outer-complexity}
\boxed{
 O\!\left(
  \Delta_P\kappa^{3/2}\sqrt\rho\,\epsilon^{-3/2}
  +\log_+\!\frac{\|\nabla P(x_0)\|}{\epsilon}+1
 \right).
}
\end{equation}
If, in addition, a module call that starts with residual $R$ uses at
most
$C_{\mathcal S}\sqrt\kappa\,[b_{\mathcal S}+
\lceil\log_2\max\{1,R/\tau_{\rm out}\}\rceil]$
inner-gradient evaluations, where $C_{\mathcal S}$ and
$b_{\mathcal S}$ are independent of $\epsilon$ and the call index,
then the total number of inner-gradient evaluations is
\begin{align}
&O\!\left(
 C_{\mathcal S}\sqrt\kappa\,b_{\mathcal S}
 \left[
  \Delta_P\sqrt{M_P}\,\epsilon^{-3/2}
  +\log_+\!\frac{\|\nabla P(x_0)\|}{\epsilon}+1
 \right]
 \right)\notag\\
&\quad+O\!\left(
 C_{\mathcal S}\sqrt\kappa
 \log_+\!\frac{
 \|\nabla_yf(x_0,y_{-1})\|}{\tau_{\rm out}}
 \right).
\label{eq:ncsc-gradient-complexity}
\end{align}
In particular, if
$C_{\mathcal S}=O(1)$ and
$b_{\mathcal S}=O(1+\log\kappa)$, its leading term is
\begin{equation}\label{eq:ncsc-gradient-leading}
 O\!\left(
  (1+\log\kappa)\Delta_P\kappa^2\sqrt\rho\,
  \epsilon^{-3/2}
 \right),
\end{equation}
with no multiplicative $\log(1/\epsilon)$.
Replacing the internal tolerance by $256\epsilon/257$ gives the
standard $(\epsilon,\sqrt{M_P\epsilon})$ guarantee without changing
these orders.
\end{theorem}




\begin{corollary}[AGD realization]
\label{cor:ncsc-agd-realization}
Choosing $\mathcal S$ to be the AGD module in
\Cref{prop:agd-implementation} satisfies the hypotheses of
\Cref{thm:ncsc-complexity} with
$C_{\mathcal S}=O(1)$ and
$b_{\mathcal S}=O(1+\log\kappa)$.  The resulting method returns the
point in \eqref{eq:ncsc-return}, has the outer complexity in
\eqref{eq:ncsc-outer-complexity}, and uses
$O((1+\log\kappa)\Delta_P\kappa^2\sqrt\rho\,\epsilon^{-3/2})$
inner-gradient evaluations in its leading term, where
$\Delta_P=P(x_0)-\inf_xP(x)$.  No
$\log(1/\epsilon)$ multiplies that term.  When $\kappa=1$, the inner
work is bounded by a numerical constant per call.
\end{corollary}

\begin{remark}[What is counted]
If the algorithm solves $K$ trust-region subproblems, it has $K$
nonterminal outer iterations and $K+1$ corrected-gradient and
Schur-complement constructions, including the terminal test.  The matrix
inversion or
linear algebra used to apply
$(\nabla^2_{yy}f)^{-1}$ is not counted as an inner gradient evaluation.
The final $\nabla_y f$ evaluation used by the residual certificate is
reused in the corrected-gradient construction.
If that inverse is implemented by Hessian-vector products, its complexity
must be reported separately.  For an accelerated implementation whose
$C_{\mathcal S}$ and $b_{\mathcal S}$ are independent of $\epsilon$,
all $\log(1/\epsilon)$ dependence in
\eqref{eq:ncsc-gradient-complexity} is additive: it occurs in the
one-time initialization term and in
$\log_+(\|\nabla P(x_0)\|/\epsilon)$.  Neither
logarithm multiplies $\Delta_P\sqrt{M_P}\epsilon^{-3/2}$.
\end{remark}

\begin{remark}[Comparison with the raw-gradient analysis]
For the raw plug-in gradient $\nabla_xf(x,y)$ used by
\citet{luo2022finding}, smoothness gives
\[
 \|\nabla_xf(x,y)-\nabla P(x)\|
 \le\ell\|y-y^*(x)\|.
\]
Thus an $O(\epsilon)$ outer-gradient error requires an
$O(\epsilon/\ell)$ inner distance.  After the summation simplification
in \citet[Theorem~2]{luo2022finding}, the $\epsilon$-dependent leading
logarithm scales, up to numerical constants, as
\[
 \log\!\left(
  \kappa^{3/2}
  +\frac{\epsilon^{3/2}}
   {\rho^{3/2}\widetilde\epsilon^3}
 \right),
 \qquad
 \widetilde\epsilon
 =\min\left\{\frac{\epsilon}{192\ell},
  \frac{\sqrt{M_P\epsilon}}{48\rho}\right\}.
\]
For sufficiently small $\epsilon$, this is
\[
 \log\!\left(
  \kappa^{3/2}
  +\left(\frac{\ell}{\sqrt{\rho\epsilon}}\right)^3
 \right).
\]
The quadratic error in \eqref{eq:corrected-gradient-error} changes the
required inner distance to order $\sqrt\epsilon$.  At the level of the
two proved upper bounds, this removes the leading multiplicative
$\log(1/\epsilon)$ while retaining the $1+\log\kappa$ dependence.
\end{remark}

\subsection{Certified adaptive specialization}

The same corrected tuple also supplies the value and derivative
certificates required by the adaptive unconstrained theorem.  We state
the oracle reduction separately so that it does not interrupt the
fixed-scale complexity argument above.

\begin{corollary}[Certified adaptive NC--SC oracle]
\label{cor:ncsc-adaptive-oracle}
Under \Cref{ass:ncsc}, let $\mathcal S$ satisfy the residual contract
\eqref{eq:maxres-interface}.  Fix $\epsilon>0$, $\omega>0$, and
$\delta_g>0$, and define
\begin{subequations}\label{eq:ncsc-adaptive-distances}
\begin{align}
 \delta_{VG}
 &:={\min}\left\{
 \sqrt{\frac{2\delta_g}{(1+\kappa)\rho}},
 \left(\frac{3\omega}{2(3\kappa-1)\rho}\right)^{1/3}
 \right\},
 \label{eq:ncsc-adaptive-vg-distance}\\
 \delta_H(\sigma)
 &:=\frac{\sigma\sqrt\epsilon}
 {64(1+\kappa)^2\rho}.
 \label{eq:ncsc-adaptive-h-distance}
\end{align}
\end{subequations}
At a value--gradient query, call $\mathcal S$ from any warm start with
target $\mu\delta_{VG}$ and return $(\widehat P,\widehat g)$.  At a
Hessian query with trial scale $\sigma$, call it with target
$\mu\delta_H(\sigma)$ and return $\widehat H$; the two queries may use
different inner points.  The returned quantities satisfy
\begin{equation}\label{eq:ncsc-adaptive-certificates}
 |\widehat P-P|\le\frac\omega8,
 \qquad
 \|\widehat g-\nabla P\|\le\delta_g,
 \qquad
 \|\widehat H-\nabla^2P\|
 \le\frac{\sigma\sqrt\epsilon}{64}.
\end{equation}
\end{corollary}

Consequently, set $M=M_P$ in
\Cref{ass:adaptive-inexact-oracle}, take $\eta=1/256$, and write
\[
 \bar\sigma:=\max\left\{\sigma_0,
  \gamma_+\sqrt{M_P/6}\right\}.
\]
The explicit choices
\begin{equation}\label{eq:ncsc-adaptive-explicit-errors}
 \omega:=\frac{\eta\epsilon^{3/2}}{16\bar\sigma},
 \qquad
 \delta_g:=\min\left\{
  \frac\epsilon{256},
  \sqrt{\frac{\omega\sqrt{M_P/6}\sqrt\epsilon}{2}}
 \right\}
\end{equation}
satisfy the compatibility conditions in that assumption.  Hence
\Cref{thm:total-complexity}, together with
\Cref{cor:ncsc-adaptive-oracle}, gives
the following bound on certified outer trials, where
$\Delta_P:=P(x_0)-P^*$:
\begin{equation}\label{eq:ncsc-adaptive-outer-complexity}
 O\!\left(
  \bar\sigma\Delta_P\epsilon^{-3/2}
  +\log_+\!\frac{\|\nabla P(x_0)\|}{\epsilon}
  +\log_+(\bar\sigma/\sigma_0)+1
 \right).
\end{equation}
The returned point satisfies
\begin{equation}\label{eq:ncsc-adaptive-return}
 \|\nabla P(\widehat x)\|
 \le\frac{257}{256}\epsilon,
 \qquad
 \lambda_{\min}(\nabla^2P(\widehat x))
 \ge-\frac{73}{64}\bar\sigma\sqrt\epsilon.
\end{equation}
This statement does not count the refinement work needed to produce
the scale-dependent certificates.

\begin{remark}[Scope of the result]
The minimax analysis uses the valid upper bound $M_P$ to fix the
trust-region scale.  Thus \eqref{eq:ncsc-outer-complexity} is a genuine
bound on all solved subproblems.  The certified adaptive theorem is
also applicable through \Cref{cor:ncsc-adaptive-oracle}; unlike the
earlier exact-oracle formulation, it explicitly accounts for corrected
value, gradient, and Hessian errors.  That corollary bounds certified
outer trials.  It does not assert the same total inner-gradient bound,
because producing its scale-dependent value certificates may require
additional refinement.  The fixed-scale result is therefore retained
for the sharp, fully accounted inner-work estimate in
\eqref{eq:ncsc-gradient-complexity}.
\end{remark}

\section{Parameter-free residual-module UTR with SCAR}
\label{sec:pf-scar-utr}

The baseline wrapper in \Cref{alg:ncsc-utr} uses the problem parameters
in both $\tau_{\rm out}$ and $M_P$; its AGD realization also uses
$\ell$ and $\mu$ internally.  Replacing AGD by SCAR instantiates the
same residual interface without inner curvature parameters, but does
not remove the parameters from the outer wrapper.  This section also
changes the outer tolerances and scale update, thereby removing both
uses.  The resulting outer wrapper takes no problem parameter as input
other than the requested tolerance $\epsilon$; with a parameter-free
residual module, the same is true of the complete method.
There is, however, an important quantifier in the theorem.  After
fixing common constants $(\ell,\mu,\rho)$, one unknown threshold
$\epsilon_0(\ell,\mu,\rho)>0$ works simultaneously for every payoff
in the corresponding class and every initialization.  The method is
therefore class-uniform asymptotically parameter-free, but the
algorithm cannot verify that $\epsilon\le\epsilon_0$.  The distinction
between this result and a uniform residual-only certificate without a
positive common lower bound on $\mu$ is detailed in
\Cref{app:scar-limitation}.

\subsection{SCAR as a residual-module realization}

We first record the precise external result needed below.  It is stated
for minimization so that it applies to
$q_x(y):=-f(x,y)$.

\begin{theorem}[SCAR; \citet{lan2026optimal}]
\label{thm:scar-blackbox}
Let $\phi:\R^m\to\R$ be differentiable, $\mu$-strongly convex, and
have an $L$-Lipschitz gradient, where $0<\mu\le L$.  Let $\tau>0$ and
$u_0\in\R^m$.  If $\|\nabla\phi(u_0)\|\le\tau$, return $u_0$.
Otherwise, initialize Algorithm~5.1 of \citet{lan2026optimal} with
$\mu_0\ge\mu$ and $0<M_0\le4L$.  SCAR returns $\widehat u$ such that
\begin{equation}\label{eq:scar-residual}
 \|\nabla\phi(\widehat u)\|\le\tau.
\end{equation}
The number of evaluations of $\nabla\phi$ is at most
\begin{equation}\label{eq:scar-original-complexity}
 (4+8\sqrt5\,C_1)\sqrt{\frac L\mu}
 \left[
  \left\lceil\log_4\frac{\mu_0}{\mu}\right\rceil
  +\left\lceil
    \log_2\frac{\|\nabla\phi(u_0)\|}{\tau}
   \right\rceil
 \right],
\end{equation}
where $C_1$ is the universal constant in Theorem~4.1 of
\citet{lan2026optimal}.  The line-search estimate returned by SCAR is
at most $2L$.
The cited theorem states the strict case $\mu<L$.  Its proof uses an
AR estimate that remains valid when $L=\mu$, so the same SCAR call and
the displayed conclusions extend to this endpoint without a
parameter-dependent branch.
\end{theorem}




SCAR can be initialized without knowing either curvature parameter.

\begin{corollary}[SCAR realization of the residual module]
\label{cor:scar-inner}
Under \Cref{ass:ncsc}, fix $x$, an inner initial point $y^{\rm in}$,
and a target $\tau>0$.  Choose any
$z\ne y^{\rm in}$ and set
\begin{equation}\label{eq:scar-secant}
 a_x:=\frac{
  \|\nabla_y f(x,z)-\nabla_y f(x,y^{\rm in})\|}
  {\|z-y^{\rm in}\|}.
\end{equation}
Run SCAR on $q_x=-f(x,\cdot)$ with
$\mu_0=M_0=a_x$ and residual tolerance $\tau$.  No value of $\ell$,
$\mu$, or $\kappa$ is supplied to the inner solver, and its output
$\widehat y$ satisfies
\[
 \|\nabla_y f(x,\widehat y)\|\le\tau.
\]
Including the secant initialization, its inner-gradient complexity is
\begin{equation}\label{eq:scar-inner-complexity}
 O\!\left(\sqrt\kappa\left[
  1+\log\kappa+
  \log_+\frac{\|\nabla_y f(x,y^{\rm in})\|}{\tau}
 \right]\right).
\end{equation}
Consequently, secant-initialized SCAR implements
\Cref{def:inner-residual-module} with
$C_{\mathcal S}=O(1)$ and
$b_{\mathcal S}=O(1+\log\kappa)$.
\end{corollary}




Thus the known-parameter outer wrapper
\Cref{alg:ncsc-utr} may use either the AGD realization in
\Cref{cor:ncsc-agd-realization} or the SCAR realization in
\Cref{cor:scar-inner}; \Cref{thm:ncsc-complexity} applies to both
without changing the outer proof.  In that baseline wrapper the target
$\tau_{\rm out}$ and scale $M_P$ are still parameter-dependent.  The
role of the present section is to change those outer quantities as
well, not merely to exchange one implementation of
$\operatorname{MaxRes}$ for another.

The published bound \eqref{eq:scar-original-complexity} counts gradient
evaluations.  SCAR is not function-value-free: its backtracking tests
also evaluate $\phi$.  With a standard implementation using a constant
number of function values per line-search trial, their number is of the
same order as \eqref{eq:scar-inner-complexity}, up to an additive
line-search logarithm.  We keep the two oracle counts conceptually
separate.

\subsection{Uniform asymptotic accuracy}

The parameter-free construction uses the corrected value
\eqref{eq:corrected-value-pf} and its error bound from
\Cref{lem:corrected-errors}.

Throughout this section, assume $0<\epsilon<1$.  Every inner call uses
the residual tolerance $\epsilon^2$, and every nonmonotone value
allowance is $\epsilon^3$.

\begin{lemma}[Class-uniform asymptotic oracle accuracy]
\label{lem:pf-asymptotic-oracle}
Fix $\ell\ge\mu>0$ and $\rho>0$, and set
$\kappa:=\ell/\mu$ and $M_P:=4\sqrt2\,\kappa^3\rho$.  There exists
$\epsilon_{\rm or}=\epsilon_{\rm or}(\ell,\mu,\rho)>0$ such that the
following holds simultaneously for every payoff satisfying
\Cref{ass:ncsc} with these common constants and every
$0<\epsilon\le\epsilon_{\rm or}$.  For every $x$ and $y$ satisfying
$\|\nabla_y f(x,y)\|\le\epsilon^2$,
\begin{subequations}\label{eq:pf-oracle-bounds}
\begin{align}
 \|\widehat g(x,y)-\nabla P(x)\|
 &\le\frac\epsilon{256},
       \label{eq:pf-gradient-error}\\
 \|\widehat H(x,y)-\nabla^2P(x)\|
 &\le\frac1{64}\sqrt{\frac{M_P\epsilon}{6}},
       \label{eq:pf-hessian-error}\\
 |\widehat P(x,y)-P(x)|
 &\le\frac{\epsilon^3}{8}.
       \label{eq:pf-value-error}
\end{align}
\end{subequations}
The same $\epsilon_{\rm or}$ works uniformly over the stated parameter
class.
\end{lemma}




\subsection{A fully observable adaptive UTR scheme}

At a current point, let $(\widehat P,g,H)$ denote the corrected tuple
computed after an inner-module call with tolerance $\epsilon^2$,
and set
$h=\max\{\|g\|,\epsilon\}$.  For a trial scale $\sigma\ge\epsilon$,
consider the padded quadratic subproblem
\begin{equation}\label{eq:pf-robust-subproblem}
\begin{aligned}
 \min_{\|d\|\le\sqrt h/(4\sigma)}\quad
 &g^Td+\frac12d^T\left(
 H+\frac{\sigma\sqrt\epsilon}{64}I
   +\sigma\sqrt h\,I
 \right)d.
\end{aligned}
\end{equation}
At every scale $\sigma^2\ge M_P/6$, the Hessian padding dominates
\eqref{eq:pf-hessian-error} and makes the quadratic term an upper
model.  For each fixed parameter class, the actual gradient error is of
the much smaller order $O(\epsilon^4)$;
at a safe scale, its worst-case effect on an interior trial is only
$O(\epsilon^{15/2})$ and is absorbed by the $O(\epsilon^3)$
nonmonotonicity.  Thus the inner solve need not be
refined to an accuracy depending on the possibly arbitrarily small
trial step.

For a global solution $d\ne0$, write $r=\|d\|$ and let $\lambda\ge0$
be the multiplier of the Euclidean ball.  Applying
\Cref{lem:tr-global-optimality} uses
$B=H+(\sigma\sqrt\epsilon/64+\sigma\sqrt h)I$ and
$R=\sqrt h/(4\sigma)$.

\begin{algorithm}[H]
\DontPrintSemicolon
\caption{Asymptotically parameter-free outer residual-module UTR}
\label{alg:pf-residual-utr}
\KwIn{$x_0,y_{-1}$, $0<\epsilon<1$, and an inner module $\mathcal S$
      satisfying \eqref{eq:maxres-interface}}
$x\leftarrow x_0$, $\sigma\leftarrow\epsilon$, and
$\eta\leftarrow1/256$\;
$y\leftarrow
 \operatorname{MaxRes}_{\mathcal S}(x_0,y_{-1},\epsilon^2)$\;
compute $(\widehat P,g,H)$ from
\eqref{eq:corrected-value-pf} and \eqref{eq:corrected-derivatives}\;
\While{true}{
 $h\leftarrow\max\{\|g\|,\epsilon\}$\;
 \If{$\|g\|\le\epsilon$ and
 $\lambda_{\min}(H)\ge-(9/8)\sigma\sqrt\epsilon$}{
  \Return{$x$}\;
 }
 Solve \eqref{eq:pf-robust-subproblem} globally and obtain
 $(d,\lambda)$; set $r\leftarrow\|d\|$\;
 \If{$r=0$}{\Return{$x$}\;}
 $x^+\leftarrow x+d$\;
 $y^+\leftarrow
  \operatorname{MaxRes}_{\mathcal S}(x^+,y,\epsilon^2)$\; compute
 $(\widehat P^+,g^+,H^+)$\;
 \eIf{
  $r<\sqrt h/(4\sigma)$,
  $\widehat P^+\le\widehat P+\epsilon^3$, and
  $\|g^+\|\le h/2$
 }{
  accept the trial\;
 }{
  \eIf{
   $r=\sqrt h/(4\sigma)$,
   $\widehat P-\widehat P^+
    \ge\eta(h^{3/2}/\sigma+\lambda r^2)+\epsilon^3$, and
   $\|g^+\|\le h/2+\lambda r$
  }{
   accept the trial\;
  }{
   reject the trial\;
  }
 }
 \eIf{the trial is accepted}{
  $(x,y,\widehat P,g,H)\leftarrow
   (x^+,y^+,\widehat P^+,g^+,H^+)$\;
  $\sigma\leftarrow\max\{\epsilon,\sigma/2\}$\;
 }{
  discard the candidate tuple and set $\sigma\leftarrow2\sigma$\;
 }
}
\end{algorithm}

Every quantity tested in \Cref{alg:pf-residual-utr} is observable.  In
particular, if $\mathcal S$ itself is parameter-free, the algorithm
never uses $\ell$, $\mu$, $\rho$, $\kappa$, $M_P$, or $\Delta_P$.
The residual target $\epsilon^2$ is independent of $\sigma$, so a
rejection changes only the robust trust-region model; the corrected
tuple at the current point remains frozen.  Secant-initialized SCAR from
\Cref{cor:scar-inner} is the implementation used for the explicit
inner-oracle count below.

\subsection{Theory of the parameter-free variant}

We next prove the trial estimate that drives the complete complexity
argument.

\begin{lemma}[Safe-scale robust trial]
\label{lem:pf-safe-trial}
Under \Cref{ass:ncsc}, fix $0<\epsilon<1$.  Consider a nonzero trial
generated by \Cref{alg:pf-residual-utr}; write
$M_P:=4\sqrt2\,\kappa^3\rho$,
$h:=\max\{\|g\|,\epsilon\}$, $r:=\|d\|$, and let $\lambda$ be its
trust-region multiplier.  Suppose the current and candidate inner
points both satisfy the module
residual bound $\|\nabla_yf\|\le\epsilon^2$, and suppose $\epsilon$ is
small enough that \eqref{eq:pf-oracle-bounds} holds.  If
\begin{equation}\label{eq:pf-safe-scale}
 \sigma^2\ge M_P/6,
\end{equation}
then a nonzero trial step satisfies
\begin{subequations}\label{eq:pf-safe-trial-bounds}
\begin{align}
 P(x)-P(x+d)
 &\ge\frac14\sigma\sqrt h\,r^2
 +\frac12\lambda r^2
 -\frac{(1+\kappa)\rho}{2\mu^2}\epsilon^4r,
       \label{eq:pf-safe-decrease}\\
 \|g^+\|
 &\le\frac{117}{256}h+\lambda r.
       \label{eq:pf-safe-gradient}
\end{align}
\end{subequations}
If $r=\sqrt h/(4\sigma)$, then also
\begin{equation}\label{eq:pf-safe-boundary-decrease}
 P(x)-P(x+d)
 \ge\frac{h^{3/2}}{128\sigma}
 +\frac14\lambda r^2
\end{equation}
without any additional smallness condition.
\end{lemma}




\begin{lemma}[Accepted-step accounting]
\label{lem:pf-accepted-accounting}
Under \Cref{ass:ncsc}, fix $0<\epsilon<1$ and consider any finite
prefix of $N$ accepted nonterminal steps generated by
\Cref{alg:pf-residual-utr}.  Assume \eqref{eq:pf-oracle-bounds} holds
at every tuple in the prefix and all its trial scales are bounded by
some $\bar\sigma>0$.  If
\begin{equation}\label{eq:pf-absorption-condition}
 \frac{5\bar\sigma\epsilon^{3/2}}{4\eta}
 \left(1+\frac4{\log2}\right)\le\frac12,
\end{equation}
then
\begin{equation}\label{eq:pf-accepted-bound}
 N=O\!\left(
 \bar\sigma\Delta_P\epsilon^{-3/2}
  +\log_+\frac{\|g_0\|}{\epsilon}+1
 \right),
\end{equation}
where $\Delta_P:=P(x_0)-P^*$, $g_0$ is the first stored corrected
gradient, and the hidden constant is numerical.  In particular, the
right-hand side is independent of the prefix length $N$.
\end{lemma}




\begin{theorem}[Class-uniform parameter-free outer NC--SC complexity]
\label{thm:pf-scar-utr}
Fix $\ell\ge\mu>0$ and $\rho>0$.  There exists
$\epsilon_0=\epsilon_0(\ell,\mu,\rho)\in(0,1)$ with the following
property.  Let $f$ be any payoff satisfying \Cref{ass:ncsc} with
these common constants and define $P(x)=\max_yf(x,y)$.  For every
initialization $(x_0,y_{-1})$, set
\[
 \kappa=\ell/\mu,
 \qquad M_P=4\sqrt2\,\kappa^3\rho,
 \qquad \Delta_P=P(x_0)-\inf_xP(x).
\]
For every $0<\epsilon\le\epsilon_0$ and every residual module $\mathcal S$
satisfying \eqref{eq:maxres-interface},
\Cref{alg:pf-residual-utr} terminates and its output $\widehat x$
satisfies
\begin{equation}\label{eq:pf-return}
 \|\nabla P(\widehat x)\|\le\frac{257}{256}\epsilon,
 \qquad
 \lambda_{\min}(\nabla^2P(\widehat x))
 \ge-\sqrt{M_P\epsilon}.
\end{equation}
Consequently, running the algorithm with internal tolerance
$256\epsilon/257$ returns an
$(\epsilon,\sqrt{M_P\epsilon})$-second-order stationary point.
The total number $Q$ of solved robust trust-region subproblems obeys
\begin{equation}\label{eq:pf-outer-complexity}
 Q=O\!\left(
  \Delta_P\sqrt{M_P}\,\epsilon^{-3/2}
  +\log_+\frac{\|\nabla P(x_0)\|}{\epsilon}
  +\log_+\frac{\sqrt{M_P}}{\epsilon}+1
 \right).
\end{equation}
Its leading problem-parameter dependence is therefore
$O(\Delta_P\kappa^{3/2}\sqrt\rho\,\epsilon^{-3/2})$.
If the module calls use none of
$\ell,\mu,\rho,\kappa,M_P$, or $\Delta_P$, then the complete
executable method uses only $\epsilon$ and universal numerical
constants.
\end{theorem}

\begin{corollary}[SCAR inner-oracle complexity]
\label{cor:pf-scar-inner-complexity}
Under the hypotheses of \Cref{thm:pf-scar-utr}, choose $\mathcal S$ to
be the secant-initialized SCAR module in \Cref{cor:scar-inner}.  If
$Q$ denotes the number of robust trust-region subproblems in
\eqref{eq:pf-outer-complexity}, set
\begin{equation}\label{eq:pf-gradient-level}
 G_Q:=1+\sqrt{2(1+\kappa)\ell
       \left(\Delta_P+\frac54Q\epsilon^3\right)}.
\end{equation}
Then the total number $N_y$ of
evaluations of $\nabla_yf$ is
\begin{align}
 N_y
 &=O\!\left(
  \sqrt\kappa\,(Q+1)
  \left[1+\log\kappa+
   \log\!\left(1+
    \frac{\ell\sqrt{G_Q}}{4\epsilon^3}
   \right)
  \right]\right)\notag\\
 &\quad+O\!\left(
  \sqrt\kappa\left[1+\log\kappa+
   \log_+\frac{\|\nabla_yf(x_0,y_{-1})\|}{\epsilon^2}
  \right]\right).
 \label{eq:pf-inner-complexity}
\end{align}
Thus the uniform SCAR solves in this construction introduce a leading
$\log(1/\epsilon)$ in the inner-oracle count, although the outer bound
in \eqref{eq:pf-outer-complexity} has no such leading factor.
\end{corollary}




\begin{remark}[Interpretation of the theorem]
The onset in \Cref{thm:pf-scar-utr} is uniform over every class with
common constants $(\ell,\mu,\rho)$ and over all initializations.  The
same executable algorithm is used without knowing these constants,
but it cannot observe whether $\epsilon\le
\epsilon_0(\ell,\mu,\rho)$.  Thus class-uniform asymptotic correctness
does not by itself provide an a posteriori finite-accuracy
certificate.
The scalar construction in
\Cref{app:scar-limitation} explains why the threshold cannot be chosen
independently of a positive common lower bound on $\mu$.

The proof gives the explicit sufficient threshold in
\eqref{eq:pf-uniform-oracle-onset} and
\eqref{eq:pf-class-uniform-onset}.  If $\ell$ and $\rho$ are fixed
while $\mu\downarrow0$, the displayed sufficient threshold is
$\Theta(\mu)$.  In
particular, a sixth-order restriction $O(\mu^6)$ is not intrinsic: it
would result from requiring the Hessian-oracle error to be smaller
than $\sigma\sqrt\epsilon$ even at the smallest trial scale
$\sigma=\epsilon$.  The proof instead needs this comparison only once
the scale is safe.  The $\epsilon^3$ value allowance is also needed
to keep the accepted-step absorption condition linear in $\mu$;
retaining an $\epsilon^2$ allowance would give a sufficient
$O(\mu^3)$ onset, which is still sharper than $O(\mu^6)$ but not
linear.
\end{remark}

\section{Adaptive fixed-radius SCAR trust region}
\label{sec:pf-fixed-radius}

This section gives a second parameter-free construction.  It combines
the fixed-radius trust-region mechanism of
\citet[Section~8.7]{luenberger2021linear} with a persistent version of
SCAR.  Persistent SCAR is a sequence-amortized realization of the same
residual interface in \Cref{def:inner-residual-module}; its cross-call
work guarantee is established separately in
\Cref{lem:scar-persistent-half}.  In contrast with
\Cref{alg:pf-residual-utr}, each ordinary inner update performs a fixed
number of observable residual halvings instead of restarting SCAR to a
prescribed power of $\epsilon$.  Consequently, neither
$\log(1/\epsilon)$ nor $\log\kappa$ multiplies the leading inner-oracle
bound.  A high-accuracy SCAR call is used only when the trust-region
solution is interior, and such validation calls occur only during the
additive scale search.

As in \Cref{thm:pf-scar-utr}, the result is class-uniform asymptotic.
For each fixed triple $(\ell,\mu,\rho)$, one unknown threshold
$\epsilon_0(\ell,\mu,\rho)>0$ works for every payoff in the class and
every initialization.  The algorithm itself uses only $\epsilon$ and
universal numerical constants; the precise quantifier and parameter
list are stated in \Cref{thm:adaptive-fixed-radius-scar}.

\subsection{Persistent SCAR and the working oracle}

We first extract from the proof of \Cref{thm:scar-blackbox} a
sequential form of SCAR.  Only the curvature calibration and the
line-search estimate are retained between calls; accelerated momentum
is restarted whenever the objective changes.

\begin{lemma}[Persistent SCAR halving]
\label{lem:scar-persistent-half}
Let $\{\phi_j\}_{j\ge1}$ be any sequence of differentiable functions
that are all $\mu$-strongly convex and have $L$-Lipschitz gradients.
Set $\kappa:=L/\mu$.
Initialize the strong-convexity guess $\nu_0$ and line-search estimate
$\mathcal M_0$ so that
\[
 \mu\le\nu_0\le L,
 \qquad 0<\mathcal M_0\le4L.
\]
A \emph{persistent halving call} at $(\phi_j,u)$ freezes
$R=\|\nabla\phi_j(u)\|$ and runs one AR call used in Algorithm~5.1 of
\citet{lan2026optimal}.  If the call returns $v$ with
$\|\nabla\phi_j(v)\|\le R/2$, the call accepts $v$ and retains the
current curvature estimates.  Otherwise it discards $v$, replaces
$\nu$ by $\nu/4$, retains the updated line-search estimate, and repeats
from $u$ on the same function.
If $R=0$, the call returns $u$ immediately and is counted as a trivial
successful halving.

The same AR call is used at the endpoint $L=\mu$.  The estimate from
Theorem~4.1 of \citet{lan2026optimal} used below remains valid there,
so this endpoint introduces no branch depending on unknown curvature
parameters.

Across any $N$ successful halving calls, the total number of failed
curvature guesses is at most
\begin{equation}\label{eq:persistent-scar-failures}
 \left\lceil\log_4\frac{\nu_0}{\mu}\right\rceil
 \le\lceil\log_4\kappa\rceil.
\end{equation}
The total number of gradient evaluations is at most
\begin{equation}\label{eq:persistent-scar-count}
 (4+2\sqrt{10}\,C_1)\sqrt\kappa
 \left(
  N+\left\lceil\log_4\frac{\nu_0}{\mu}\right\rceil
 \right),
\end{equation}
where $C_1$ is the constant in Theorem~4.1 of
\citet{lan2026optimal}, also used in
\eqref{eq:scar-original-complexity}.  After any nontrivial AR call, the
retained line-search estimate is at most $2L$; before the first such
call, it remains at its admissible initial value at most $4L$.
\end{lemma}




For the rest of this section, define
\begin{equation}\label{eq:fixed-radius-parameters}
 \mathcal A_\epsilon:=\log(e/\epsilon),
 \qquad
 \sigma_0:=\mathcal A_\epsilon^{-1},
 \qquad
 c_0:=2^{-12},
 \qquad
 R(\sigma):=\frac{\sqrt\epsilon}{4\sigma},
\end{equation}
and let
\begin{equation}\label{eq:working-half-count}
 N_{\rm w}:=
 \left\lceil\log_2\frac{1+c_0}{c_0}\right\rceil.
\end{equation}
A corrected tuple $(\widehat P,\widehat g,\widehat H)$ at $(x,y)$ and
scale $\sigma$ is called \emph{working-accurate} if
\begin{equation}\label{eq:working-residual-invariant}
 \|\nabla_yf(x,y)\|
 \le c_0\ell R(\sigma).
\end{equation}
This condition is used only in the analysis; the algorithm maintains
it by a fixed number of observable residual halvings and never uses
$\ell$.

\begin{lemma}[Working residual invariant]
\label{lem:working-residual-invariant}
Under \Cref{ass:ncsc}, fix $0<\epsilon<1$ and $\sigma>0$.  Suppose
\eqref{eq:working-residual-invariant} holds at $(x,y)$.
If $\|d\|\le R(\sigma)$, then $N_{\rm w}$ persistent halving calls on
$q_{x+d}=-f(x+d,\cdot)$, starting from $y$, produce $y^+$ satisfying
\[
 \|\nabla_yf(x+d,y^+)\|
 \le c_0\ell R(\sigma).
\]
If the scale is doubled at the same $x$, one persistent halving call
produces a point satisfying
\eqref{eq:working-residual-invariant} with scale $2\sigma$.
\end{lemma}




The invariant implies
\begin{equation}\label{eq:working-distance-bound}
 \|y-y^*(x)\|
 \le c_0\kappa R(\sigma)
 =\frac{c_0\kappa\sqrt\epsilon}{4\sigma}.
\end{equation}
For later accounting, introduce the analytical, nonobservable error
majorant
\begin{equation}\label{eq:fixed-radius-proxy-error}
 \mathcal E_\epsilon(\sigma):=
 \frac{7\rho}{24}\bigl(c_0\kappa R(\sigma)\bigr)^3
 =\frac{7c_0^3\rho\kappa^3}{1536}
       \frac{\epsilon^{3/2}}{\sigma^3}.
\end{equation}

\begin{lemma}[Class-uniform working and validation accuracy]
\label{lem:fixed-radius-oracle-accuracy}
Fix $\ell\ge\mu>0$ and $\rho>0$, let
$\kappa=\ell/\mu$ and $M_P=4\sqrt2\kappa^3\rho$, and consider any
payoff satisfying \Cref{ass:ncsc} with these common constants.  There
exist $\epsilon_1=\epsilon_1(\ell,\mu,\rho)>0$ and
$\epsilon_2=\epsilon_2(\ell,\mu,\rho)>0$ such that the following
statements hold uniformly over this parameter class.
\begin{enumerate}[label=(\roman*)]
 \item For every $0<\epsilon\le\epsilon_1$, every
 $\sigma\ge\mathcal A_\epsilon^{-1}$, and every working-accurate tuple,
 \begin{equation}\label{eq:working-proxy-bound}
  |\widehat P(x,y)-P(x)|
  \le\mathcal E_\epsilon(\sigma).
 \end{equation}
 \item After decreasing $\epsilon_1$ so that
 $\mathcal A_\epsilon^{-1}\le\sqrt{M_P/6}$ for all
 $0<\epsilon\le\epsilon_1$, every working-accurate tuple with
 $\sigma^2\ge M_P/6$ satisfies
 \begin{equation}\label{eq:working-safe-errors}
 \begin{aligned}
  \|\widehat g-\nabla P(x)\|&\le\epsilon/1024,\\
  \|\widehat H-\nabla^2P(x)\|&\le\sigma\sqrt\epsilon/64,\\
  |\widehat P-P(x)|&\le\epsilon^{3/2}/(4096\sigma).
 \end{aligned}
 \end{equation}
 \item Choose any secant $a_x$ as in \eqref{eq:scar-secant}.  If a
 SCAR validation call at $x$ returns $y_{\rm v}$ satisfying
 \begin{equation}\label{eq:validator-residual}
  \|\nabla_yf(x,y_{\rm v})\|
  \le c_0a_x\epsilon^{3/2},
 \end{equation}
 then, for all $0<\epsilon\le\epsilon_2$ and all
 $\sigma\ge\mathcal A_\epsilon^{-1}$,
 \begin{equation}\label{eq:validator-errors}
 \begin{aligned}
  \|\widehat g(x,y_{\rm v})-\nabla P(x)\|
   &\le\epsilon/1024,\\
  \|\widehat H(x,y_{\rm v})-\nabla^2P(x)\|
   &\le\sigma\sqrt\epsilon/32.
 \end{aligned}
 \end{equation}
\end{enumerate}
The same $\epsilon_1$ and $\epsilon_2$ work independently of the
payoff, $x$, the tuple, and the chosen secant within the stated class.
\end{lemma}




\subsection{The adaptive fixed-radius algorithm}

Given a working-accurate corrected tuple, solve
\begin{equation}\label{eq:adaptive-fixed-radius-subproblem}
\begin{aligned}
 \min_{\|d\|\le R(\sigma)}\quad
 &\widehat g^Td+\frac12d^T
 \left(\widehat H+\frac{65}{64}\sigma\sqrt\epsilon I\right)d.
\end{aligned}
\end{equation}
The radius is fixed for a given $(\epsilon,\sigma)$ and does not depend
on $\|\widehat g\|$.  If $d$ is a global solution and $\lambda\ge0$
is its trust-region multiplier, then
\Cref{lem:tr-global-optimality} applies with
$B=\widehat H+(65/64)\sigma\sqrt\epsilon I$ and
$R=R(\sigma)$.

\begin{algorithm}[H]
\DontPrintSemicolon
\caption{Parameter-free adaptive fixed-radius TR with persistent SCAR}
\label{alg:adaptive-fixed-radius-scar}
\KwIn{$x_0,y_{-1}$ and $0<\epsilon<1$}
$\mathcal A_\epsilon\leftarrow\log(e/\epsilon)$,
$\sigma\leftarrow\mathcal A_\epsilon^{-1}$,
$c_0\leftarrow2^{-12}$, and $x\leftarrow x_0$\;
Choose $z_0\ne y_{-1}$, compute the secant $a_0$ from
\eqref{eq:scar-secant}, and initialize the persistent SCAR state with
$\nu_0=\mathcal M_0=a_0$\;
At $x_0$, apply persistent SCAR halving from $y_{-1}$ until
$\|\nabla_yf(x_0,y)\|\le c_0a_0R(\sigma)$; compute the corrected tuple
$(\widehat P,\widehat g,\widehat H)$\;
\While{true}{
 Solve \eqref{eq:adaptive-fixed-radius-subproblem} globally and obtain
 $d$; set $x^+\leftarrow x+d$\;
 Starting from $y$, apply $N_{\rm w}$ successful persistent SCAR
 halvings to $q_{x^+}$ and call the result $y^+$\;
 \eIf{$\|d\|=R(\sigma)$}{
  compute $\widehat P^+$ at $(x^+,y^+)$\;
  \eIf{$\widehat P-\widehat P^+
        \ge\epsilon^{3/2}/(256\sigma)$}{
   compute $(\widehat g^+,\widehat H^+)$ at $(x^+,y^+)$\;
   $(x,y,\widehat P,\widehat g,\widehat H)
    \leftarrow(x^+,y^+,\widehat P^+,\widehat g^+,\widehat H^+)$\;
  }{
   discard the candidate tuple, set $\sigma\leftarrow2\sigma$,
   apply one successful persistent halving call at the unchanged $x$,
   replace $y$ by the returned inner point, and recompute the current
   corrected tuple\;
  }
 }{
  Choose $z^+\ne y^+$ and compute at $x^+$ the secant $a^+$ from
  \eqref{eq:scar-secant}\;
  Continue persistent SCAR halving at $x^+$ until
  $\|\nabla_yf(x^+,y_{\rm v})\|
    \le c_0a^+\epsilon^{3/2}$; compute
  $(\widehat g_{\rm v},\widehat H_{\rm v})$\;
  \eIf{$\|\widehat g_{\rm v}\|\le\epsilon/2$ and
   $\lambda_{\min}(\widehat H_{\rm v})
    \ge-(21/8)\sigma\sqrt\epsilon$}{
   \Return{$x^+$}\;
  }{
   discard the candidate and validator tuples, set
   $\sigma\leftarrow2\sigma$, apply one successful persistent halving
   call at the unchanged $x$, replace $y$ by the returned inner point,
   and recompute the current corrected tuple\;
  }
 }
}
\end{algorithm}

Every operation in \Cref{alg:adaptive-fixed-radius-scar} is observable.
The secants are used only to initialize SCAR and to specify the
high-accuracy validation residual.  Strong monotonicity and smoothness
give $\mu\le a_0,a^+\le\ell$, but neither endpoint is known to the
algorithm.  After an accepted boundary step, the candidate corrected
tuple is formed at the already computed inner point and stored without
further refinement.  After a failed trial, the candidate is discarded,
the current point is unchanged, and the refreshed tuple at
the doubled scale replaces the stored tuple.  This bookkeeping is
essential for the proxy telescope below.  The scale never decreases.
The persistent SCAR state $(\nu,\mathcal M)$ is retained after every
accepted trial, rejected trial, and failed validation; rejecting an
outer candidate never rolls back this state.  The initial residual
evaluation is counted as the first gradient evaluation of the
corresponding AR call, and the final residual gradient is reused when
forming the corrected tuple.

The initialization satisfies the working invariant because
$a_0\le\ell$.  Thereafter, \Cref{lem:working-residual-invariant} proves
inductively that every tuple used in
\eqref{eq:adaptive-fixed-radius-subproblem} is working-accurate.

\subsection{Safe scale, termination, and outer complexity}

The following lemma is the fixed-radius analogue of the large-scale
acceptance result for adaptive UTR.

\begin{lemma}[Class-uniform safe fixed-radius trial]
\label{lem:adaptive-fixed-radius-safe}
Fix $\ell\ge\mu>0$ and $\rho>0$, set
$\kappa=\ell/\mu$ and $M_P:=4\sqrt2\,\kappa^3\rho$.  There exists
$\epsilon_3=\epsilon_3(\ell,\mu,\rho)>0$ such that the following
holds simultaneously for every payoff satisfying \Cref{ass:ncsc}
with these common constants and every
$0<\epsilon\le\epsilon_3$.  Consider any trial generated by
\Cref{alg:adaptive-fixed-radius-scar}, and suppose its current and
candidate tuples are working-accurate and its validator, when invoked,
satisfies \eqref{eq:validator-residual}.  Let $d$ be its global
trust-region solution and let $\lambda$ be the multiplier supplied by
\Cref{lem:tr-global-optimality} with
$B=\widehat H+(65/64)\sigma\sqrt\epsilon I$ and $R=R(\sigma)$.
Assertions (i)--(ii) below hold when
the trial scale satisfies
\begin{equation}\label{eq:adaptive-fixed-radius-safe-scale}
 \sigma^2\ge M_P/6,
\end{equation}
and assertion (iii) holds at every generated scale
$\sigma\ge\mathcal A_\epsilon^{-1}$.
\begin{enumerate}[label=(\roman*)]
 \item If $\|d\|=R(\sigma)$, then
 \begin{equation}\label{eq:safe-boundary-true-decrease}
  P(x)-P(x+d)
  \ge\frac{63}{4096}\frac{\epsilon^{3/2}}{\sigma}
       +\frac\lambda2R(\sigma)^2,
 \end{equation}
 and the observable boundary acceptance test succeeds.
 \item If $\|d\|<R(\sigma)$, then, at $x^+=x+d$,
 \begin{equation}\label{eq:safe-interior-true-bounds}
 \begin{aligned}
  \|\nabla P(x^+)\|&\le457\epsilon/1024,\\
  \lambda_{\min}(\nabla^2P(x^+))
   &\ge-(81/32)\sigma\sqrt\epsilon.
 \end{aligned}
 \end{equation}
 The validation test succeeds.
 \item At any scale, not necessarily a safe one, success of the
 validation test implies
 \begin{equation}\label{eq:validator-true-return}
  \|\nabla P(x^+)\|<\epsilon,
  \qquad
  \lambda_{\min}(\nabla^2P(x^+))
  \ge-(85/32)\sigma\sqrt\epsilon.
 \end{equation}
\end{enumerate}
\end{lemma}




\begin{theorem}[Class-uniform parameter-free fixed-radius complexity]
\label{thm:adaptive-fixed-radius-scar}
Fix $\ell\ge\mu>0$ and $\rho>0$.  There exists
$\epsilon_0=\epsilon_0(\ell,\mu,\rho)\in(0,1)$ with the following
property.  Let $f$ be any payoff satisfying \Cref{ass:ncsc} with
these common constants, define $P(x)=\max_y f(x,y)$, and fix any
initialization $(x_0,y_{-1})$.  Set
\[
 \kappa=\ell/\mu,
 \qquad M_P=4\sqrt2\,\kappa^3\rho,
 \qquad \Delta_P=P(x_0)-\inf_xP(x).
\]
Then, for every $0<\epsilon\le\epsilon_0$,
\Cref{alg:adaptive-fixed-radius-scar} terminates.  The algorithm uses
only $\epsilon$ and universal numerical constants; in particular, it
does not use $\ell,\mu,\rho,\kappa,M_P$, or $\Delta_P$.  Its returned
point $x_{\rm out}$ satisfies
\begin{equation}\label{eq:fixed-radius-return}
 \|\nabla P(x_{\rm out})\|<\epsilon,
 \qquad
 \lambda_{\min}(\nabla^2P(x_{\rm out}))
 \ge-\frac{85}{16\sqrt6}\sqrt{M_P\epsilon}.
\end{equation}
If $Q$ is the total number of solved trust-region subproblems and
$N_y$ is the total number of evaluations of $\nabla_yf$, including
the secant evaluations, then
\begin{subequations}\label{eq:fixed-radius-complexities}
\begin{align}
 Q
 &=O\!\left(\Delta_P\sqrt{M_P}\,\epsilon^{-3/2}\right)
   +O_{\rm fixed}\!\left(\operatorname{polylog}(1/\epsilon)\right),
 \label{eq:fixed-radius-leading-outer}\\
 N_y&=
 O\!\left(
  \Delta_P\sqrt{\kappa M_P}\,\epsilon^{-3/2}
 \right)
 +O_{\rm fixed}\!\left(
  \sqrt\kappa\,\operatorname{polylog}(1/\epsilon)
 \right).
 \label{eq:fixed-radius-inner-leading}
\end{align}
\end{subequations}
Here $O_{\rm fixed}$ may depend on the fixed payoff and initialization
but not on $\epsilon$.  Equivalently, the leading outer and inner
terms are, respectively,
\[
 O(\Delta_P\kappa^{3/2}\sqrt\rho\,\epsilon^{-3/2})
 \quad\text{and}\quad
 O(\Delta_P\kappa^2\sqrt\rho\,\epsilon^{-3/2}).
\]
Thus neither $\log(1/\epsilon)$ nor $\log\kappa$ multiplies the leading
inner-gradient term.  Running the method with internal tolerance
$(1536/7225)\epsilon$ gives the standard
$(\epsilon,\sqrt{M_P\epsilon})$ guarantee, provided that internal
tolerance lies in the class-uniform regime above.
\end{theorem}

\begin{remark}[Scope of the fixed-radius guarantee]
The onset in \Cref{thm:adaptive-fixed-radius-scar} is uniform over the
class with common constants $(\ell,\mu,\rho)$ and over all
initializations, but it is not observable to the parameter-free
algorithm.  The residual-only obstruction for prescribed absolute
tolerances, its precise relation to the secant-scaled targets used here,
and the roles of the corrected-value telescope and the interior
validator are explained in \Cref{app:fixed-radius-scope}.
\end{remark}




\begin{remark}[Oracle accounting]
The bound in \Cref{thm:adaptive-fixed-radius-scar} counts
inner-gradient evaluations, which is the oracle counted in
\Cref{thm:scar-blackbox}.  SCAR's line search also evaluates
$q_x=-f(x,\cdot)$.  Under the standard implementation in which the
accelerated subroutine uses $O(1)$ function values per gradient
evaluation, the function-evaluation remark following Theorem~4.1 of
\citet{lan2026optimal} bounds the extra cost of an AR call containing
$S_{\rm AR}$ backtracking stages by
$2S_{\rm AR}+\log_2(\mathcal M^+/\mathcal M)$.  Summed over all
persistent calls, the stage counts are absorbed by $N_y$, while the
logarithmic ratios telescope because $\mathcal M$ is persistent.  In
particular,
\[
 \log_2(\mathcal M_{\rm final}/a_0)
 \le\log_2(2\ell/a_0)\le1+\log_2\kappa.
\]
Thus the total inner function-value count is
$O(N_y+\log\kappa)$.  The initial corrected tuple, trial candidates,
and scale refreshes use at most $2Q+1$ additional payoff values.
Each stored tuple and each validator also constructs corrected
gradient and Hessian estimators.  The inverse actions in
\eqref{eq:corrected-value-pf} and
\eqref{eq:corrected-derivatives} are treated as exact linear algebra,
as in the earlier minimax analysis; iterative linear solves require a
separate accuracy count.
\end{remark}

\section{Conclusion}

For unconstrained optimization, the certified-inexact adaptive UTR
removes $M$ from the outer scale update and admits the value and
derivative certificates needed by the minimax reduction.  Producing
nonzero certified radii remains an oracle-interface requirement.  The
exact-oracle specialization is fully parameter free and retains sharp
leading dependence on $\Delta_F$, $M$, and $\epsilon$ with only
additive search overhead.
The fixed-scale inexact analysis is used as an outer estimate rather
than a separate algorithm; in particular, its
cubic-path budget is what amortizes the inner work in the minimax
specialization.

For NC--SC minimax optimization, the outer method now receives a
residual-based maximization module as an input.  Correctness depends
only on the module's returned residual, while AGD, restarted SCAR, and
persistent SCAR provide different work guarantees for the same
interface.  The Newton-corrected value-function gradient permits inner
solves at a $\sqrt\epsilon$ scale.  This preserves the
$O(\Delta_P\kappa^{3/2}\sqrt\rho\,\epsilon^{-3/2})$ outer complexity
and removes the leading multiplicative $\log(1/\epsilon)$ from the
raw-gradient tracking bound, while retaining logarithmic dependence on
$\kappa$ for the restarted accelerated implementations.

The first parameter-free residual-module UTR removes all problem
parameters from the executable outer and inner algorithms and retains
the sharp outer leading order uniformly over each class with common
constants $(\ell,\mu,\rho)$ once $\epsilon$ is sufficiently small,
but its uniformly accurate inner solves restore a leading
$\log(1/\epsilon)$.  The adaptive fixed-radius
variant replaces those ordinary solves by persistent residual-halving
calls.  It preserves the same leading outer order and attains leading
inner-gradient complexity
$O(\Delta_P\kappa^2\sqrt\rho\,\epsilon^{-3/2})$, with all
$\log(1/\epsilon)$ and $\log\kappa$ terms confined to additive
initialization, validation, and scale-search costs.  Both
parameter-free conclusions are asymptotic but class-uniform: their
onset depends only on the common regularity constants and not on the
payoff or initialization.  For the first construction, the displayed
sufficient onset is $\Theta(\mu)$ for fixed $\ell$ and $\rho$.  Because
the executable algorithms neither know nor certify the relevant onset,
this is still weaker than an a posteriori finite-accuracy guarantee.

\appendix
\section{Deferred technical proofs}
\label{app:proofs}
The main text states the assumptions and results used by subsequent sections.  This appendix collects the technical proofs in their order of appearance.

\subsection{A common inexact trust-region estimate}
\label{app:common-tr-estimate}

\begin{lemma}[Common inexact trust-region step estimate]
\label{lem:common-inexact-tr-step}
Suppose \Cref{ass:generic} holds.  At a point $x$, let $g$ and a
symmetric $H$ satisfy
\[
 \|g-\nabla F(x)\|\le\delta_0,
 \qquad
 \|H-\nabla^2F(x)\|\le\xi.
\]
For $a\ge0$ and $R>0$, let $d$ be a global solution of
\[
 \min_{\|u\|\le R}\;
 q(u):=g^Tu+\frac12u^T(H+aI)u,
\]
let $\lambda$ be the multiplier supplied by
\Cref{lem:tr-global-optimality}, and write $r=\|d\|$.  Let a candidate
gradient $g^+$ satisfy
$\|g^+-\nabla F(x+d)\|\le\delta_+$.  Then
\begin{subequations}\label{eq:common-tr-estimates}
\begin{align}
 q(d)&\le-\frac{\lambda}{2}r^2,
       \label{eq:common-model-bound}\\
 F(x)-F(x+d)
 &\ge\frac{a-\xi}{2}r^2+\frac{\lambda}{2}r^2
      -\frac M6r^3-\delta_0r,
       \label{eq:common-function-bound}\\
 \|g^+\|
 &\le(a+\xi+\lambda)r+\frac M2r^2+\delta_0+\delta_+.
       \label{eq:common-gradient-bound}
\end{align}
\end{subequations}
\end{lemma}

\begin{proof}
Apply \eqref{eq:tr-global-optimality} with $B=H+aI$.  Its stationarity
and positive-semidefiniteness conditions give
\[
 q(d)=-\frac12d^T(H+aI)d-\lambda r^2
 \le-\frac{\lambda}{2}r^2,
\]
which proves \eqref{eq:common-model-bound}.  The Taylor bound and the
current endpoint errors imply
\[
 F(x)-F(x+d)
 \ge-g^Td-\frac12d^THd-\delta_0r-\frac{\xi}{2}r^2-\frac M6r^3.
\]
Since
$-g^Td-d^THd/2=-q(d)+ar^2/2$, the model bound yields
\eqref{eq:common-function-bound}.

For the candidate gradient, decompose the two endpoint errors, Taylor
remainder, Hessian error, and model residual:
\begin{align*}
g^+
&=[g^+-\nabla F(x+d)]
 +[\nabla F(x+d)-\nabla F(x)-\nabla^2F(x)d]\\
&\quad+[\nabla F(x)-g]
 +[\nabla^2F(x)-H]d+[g+Hd].
\end{align*}
The stationarity part of \eqref{eq:tr-global-optimality}, again with
$B=H+aI$, gives $g+Hd=-(a+\lambda)d$.  Taking norms proves
\eqref{eq:common-gradient-bound}.
\end{proof}

\subsection{Fixed-scale inexact UTR estimates}
\label{app:fixed-scale-proofs}

Let $K\in\mathbb N_0\cup\{\infty\}$ be the first return index of the
procedure \eqref{eq:generic-stopping-test}--\eqref{eq:generic-subproblem},
with $K=\infty$ if it never returns.  For $k<K$, set
$r_k:=\|d_k\|$ and let $\lambda_k$ be the multiplier supplied by
\Cref{lem:tr-global-optimality} with
\[
 g=g_k,\qquad
 B=H_k+\sqrt{M/6}\sqrt{h_k}\,I,\qquad
 R=\frac{\sqrt{h_k}}{4\sqrt{M/6}}.
\]

\begin{lemma}[Decrease and clipped-gradient recurrence]
\label{lem:one-step}
Fix $\epsilon>0$ and suppose
\Cref{ass:generic,ass:inexact-oracle} hold.  Every nonterminal step
$k<K$ generated by \eqref{eq:generic-stopping-test}--
\eqref{eq:generic-subproblem} satisfies
\begin{equation}\label{eq:basic-decrease}
 F(x_k)-F(x_{k+1})
 \ge \frac{31}{128}\sqrt{\frac M6}\sqrt{h_k}\,r_k^2
     +\frac12\lambda_kr_k^2-\frac{\epsilon r_k}{256}.
\end{equation}
Consequently,
\begin{subequations}\label{eq:step-type-decrease}
\begin{align}
 \lambda_k=0:\quad
 F(x_k)-F(x_{k+1})
 &\ge-\frac{\epsilon^{3/2}}{63488\sqrt{M/6}},
      \label{eq:interior-small-increase}\\
 \lambda_k>0:\quad
 F(x_k)-F(x_{k+1})
 &\ge\frac1{128}
  \left(\frac{h_k^{3/2}}{\sqrt{M/6}}+\lambda_kr_k^2\right).
      \label{eq:boundary-decrease}
\end{align}
\end{subequations}
Moreover,
\begin{equation}\label{eq:gradient-recurrence}
 h_{k+1}\le
 \max\left\{\frac{15}{32}h_k+\lambda_kr_k,\epsilon\right\}.
\end{equation}
\end{lemma}

The small negative term in \eqref{eq:interior-small-increase} is why a
uniform $O(\epsilon)$ gradient error is admissible.  An inexact step
with $\lambda_k=0$ need not be monotone: it may increase $F$ by at
most $\epsilon^{3/2}/(63488\sqrt{M/6})$.  The finite-prefix argument
below debits every such increase against the decrease on indices in
$\mathcal B_N$.

\begin{proof}[Proof of \Cref{lem:one-step}]
Apply \Cref{lem:common-inexact-tr-step} with
\[
 a=\sqrt{M/6}\sqrt{h_k},\qquad
 R=\frac{\sqrt{h_k}}{4\sqrt{M/6}},\qquad
 \xi=\frac{\sqrt{M/6}\sqrt\epsilon}{64},\qquad
 \delta_0=\delta_+=\frac{\epsilon}{256}.
\]
Because $r_k\le R$,
$(M/6)r_k^3\le\sqrt{M/6}\sqrt{h_k}\,r_k^2/4$.
Using $\sqrt\epsilon\le\sqrt{h_k}$ in
\eqref{eq:common-function-bound} therefore gives
\[
 F(x_k)-F(x_{k+1})
 \ge\frac{31}{128}\sqrt{\frac M6}\sqrt{h_k}\,r_k^2
  +\frac12\lambda_kr_k^2-\frac{\epsilon r_k}{256},
\]
which is \eqref{eq:basic-decrease}.

If $\lambda_k=0$, minimizing
$31\sqrt{M/6}\sqrt\epsilon\,r^2/128-\epsilon r/256$ over $r\ge0$
has minimizer $r=\sqrt\epsilon/(124\sqrt{M/6})$ and minimum
$-\epsilon^{3/2}/(63488\sqrt{M/6})$, proving
\eqref{eq:interior-small-increase}.  If $\lambda_k>0$, then
$r_k=\sqrt{h_k}/(4\sqrt{M/6})$, and
\[
  \frac{31}{128}\sqrt{\frac M6}\sqrt{h_k}\,r_k^2
  -\frac{\epsilon r_k}{256}
  \ge\left(\frac{31}{2048}-\frac1{1024}\right)
      \frac{h_k^{3/2}}{\sqrt{M/6}}
  =\frac{29}{2048}\frac{h_k^{3/2}}{\sqrt{M/6}}.
\]
Since $29/2048>1/128$, the boundary estimate follows.

The common gradient estimate gives
\begin{align*}
\|g_{k+1}\|
&\le \sqrt{\frac M6}\sqrt{h_k}\,r_k+\lambda_kr_k
  +\frac M2r_k^2+\frac{\sqrt{M/6}\sqrt\epsilon}{64}r_k
  +\frac{\epsilon}{128}\\
&\le\left(\frac14+\frac3{16}+\frac1{256}
  +\frac2{256}\right)h_k+\lambda_kr_k\\
&=\frac{115}{256}h_k+\lambda_kr_k
 <\frac{15}{32}h_k+\lambda_kr_k.
\end{align*}
Here we used, in order,
$\sqrt{M/6}\sqrt{h_k}r_k\le h_k/4$,
$(M/2)r_k^2\le3h_k/16$,
$\sqrt{M/6}\sqrt\epsilon r_k/64\le h_k/256$, and
$\epsilon/128\le2h_k/256$.
Clipping the resulting bound below by $\epsilon$ proves
\eqref{eq:gradient-recurrence}.
\end{proof}

For any finite prefix of $N\le K$ nonterminal iterations, let
\[
 \mathcal B_N:=\{0\le k<N:\lambda_k>0\},
 \qquad
 \mathcal I_N:=\{0\le k<N:\lambda_k=0\},
\]
and, for $k\in\mathcal B_N$, define
\begin{equation}\label{eq:z-definition}
 z_k:=\frac{\lambda_kr_k}{h_k}.
\end{equation}

\begin{lemma}[Boundary and interior budgets]
\label{lem:budgets}
Fix $\epsilon>0$ and suppose
\Cref{ass:generic,ass:inexact-oracle} hold.  For the sequence generated
by \eqref{eq:generic-stopping-test}--\eqref{eq:generic-subproblem},
every finite nonterminal prefix satisfies
\begin{subequations}\label{eq:boundary-budgets}
\begin{align}
 |\mathcal B_N|
 &\le128\sqrt{\frac M6}\Delta_F\epsilon^{-3/2}
      +\frac1{496}|\mathcal I_N|,
      \label{eq:boundary-count}\\
 \sum_{k\in\mathcal B_N}z_k
 &\le512\sqrt{\frac M6}\Delta_F\epsilon^{-3/2}
      +\frac1{124}|\mathcal I_N|.
      \label{eq:z-budget}
\end{align}
\end{subequations}
Furthermore,
\begin{equation}\label{eq:interior-budget}
 |\mathcal I_N|
 \le |\mathcal B_N|+1+
 \frac{\log_+(h_0/\epsilon)+\sum_{k\in\mathcal B_N}z_k}
      {\log(32/15)}.
\end{equation}
\end{lemma}

\begin{proof}[Proof of \Cref{lem:budgets}]
Summing \eqref{eq:step-type-decrease} and using
$F(x_N)\ge F^*$ gives
\[
  \Delta_F
  \ge\frac{\epsilon^{3/2}}{128\sqrt{M/6}}|\mathcal B_N|
    +\frac1{128}\sum_{k\in\mathcal B_N}\lambda_kr_k^2
    -\frac{\epsilon^{3/2}}{63488\sqrt{M/6}}|\mathcal I_N|.
\]
This proves \eqref{eq:boundary-count}.  For $k\in\mathcal B_N$,
\[
  \lambda_kr_k^2
  =\frac{z_kh_k^{3/2}}{4\sqrt{M/6}}
  \ge\frac{z_k\epsilon^{3/2}}{4\sqrt{M/6}},
\]
which proves \eqref{eq:z-budget}.

For $k\in\mathcal B_N$, \eqref{eq:gradient-recurrence} implies
\[
 \left[\log\frac{h_{k+1}}{h_k}\right]_+
 \le\left[\log\left(\frac{15}{32}+z_k\right)\right]_+
 \le z_k.
\]
Indeed, the left-hand side is zero when $z_k\le17/32$; above this
threshold, $z-\log(15/32+z)$ is increasing and positive.
For $k\in\mathcal I_N$ with $h_{k+1}>\epsilon$, the quantity
$\log h_k$ decreases by at least $\log(32/15)$.  If
$k\in\mathcal I_N$ has
$h_{k+1}=\epsilon$ and $k+1<N$, the gradient stopping test holds at the
next iterate.  Since the repeated-subproblem procedure has not
terminated,
$\lambda_{\min}(H_{k+1})<-\sqrt{M/6}\sqrt\epsilon$, so
$H_{k+1}+\sqrt{M/6}\sqrt\epsilon I$ is not positive semidefinite.
Applying \eqref{eq:tr-global-optimality} at iteration $k+1$ with
$B=H_{k+1}+\sqrt{M/6}\sqrt\epsilon I$ therefore forces
$\lambda_{k+1}>0$, so the next index belongs to $\mathcal B_N$.
These floor-hitting indices in
$\mathcal I_N$ pair injectively with their successors in
$\mathcal B_N$, apart from the possible last index $k=N-1$.  Hence
\[
 \bigl|\{k\in\mathcal I_N:h_{k+1}=\epsilon\}\bigr|
 \le|\mathcal B_N|+1.
\]
For the remaining indices in $\mathcal I_N$, telescoping $\log h_k$,
using $h_N\ge\epsilon$, and charging every positive increase on
$k\in\mathcal B_N$ to $z_k$ gives
\[
 \bigl|\{k\in\mathcal I_N:h_{k+1}>\epsilon\}\bigr|
 \log(32/15)
 \le\log_+(h_0/\epsilon)+\sum_{k\in\mathcal B_N}z_k.
\]
Adding the last two displays proves \eqref{eq:interior-budget}.
\end{proof}

\begin{proof}[Proof of \Cref{thm:generic-complexity}]
Substituting \eqref{eq:boundary-budgets} into
\eqref{eq:interior-budget} gives
\begin{align*}
&\left[1-\frac1{496}
  -\frac1{124\log(32/15)}\right]
  |\mathcal I_N|\\
&\quad\le
128\left(1+\frac4{\log(32/15)}\right)
\sqrt{\frac M6}\Delta_F\epsilon^{-3/2}
+1+\frac{\log_+(h_0/\epsilon)}{\log(32/15)}.
\end{align*}
The coefficient in brackets is greater than $0.98$.  Hence both
$|\mathcal I_N|$ and $|\mathcal B_N|$ are bounded by
\[
 O\!\left(
 \sqrt{M/6}\,\Delta_F\epsilon^{-3/2}
 +\log_+(h_0/\epsilon)+1
 \right),
\]
uniformly over every finite prefix.
If $K=\infty$, then for every finite $N$ the identity
$N=|\mathcal I_N|+|\mathcal B_N|$ would be bounded by the same constant,
independent of $N$.  Choosing a larger $N$ gives a contradiction.
Therefore $K<\infty$.  Taking $N=K$ proves
\eqref{eq:generic-complexity} for the actual
number of solved subproblems.

At termination, \eqref{eq:generic-oracle-errors} gives
$\|\nabla F(\widehat x)\|\le257\epsilon/256$ and
\[
  \lambda_{\min}(\nabla^2F(\widehat x))
  \ge-\frac{65}{64}\sqrt{\frac M6}\sqrt\epsilon
  \ge-\sqrt{M\epsilon},
\]
where the last inequality uses $65/(64\sqrt6)<1$.
\end{proof}

\begin{proof}[Proof of \Cref{lem:cubic-path}]
Since $\sqrt{h_k}\ge4\sqrt{M/6}\,r_k$,
\eqref{eq:basic-decrease} implies
\[
 F(x_k)-F(x_{k+1})
 \ge\frac{31}{192}Mr_k^3-\frac{\epsilon r_k}{256}.
\]
Split the first term equally and maximize
$\epsilon r/256-(31/384)Mr^3$ over $r\ge0$.  Its maximum is
$\epsilon^{3/2}/(384\sqrt{62M})$, and therefore
\[
 F(x_k)-F(x_{k+1})
 \ge\frac{31}{384}Mr_k^3
 -\frac{\epsilon^{3/2}}{384\sqrt{62M}}.
\]
Summing from $k=0$ to $K-1$ and using $F(x_K)\ge F^*$ proves
\eqref{eq:cubic-path}.
\end{proof}

\subsection{Certified-inexact adaptive clipped UTR}

\begin{lemma}[Guaranteed trial outcome]\label{lem:large-scale}
Suppose \Cref{ass:adaptive-inexact-oracle} holds.  If
\begin{equation}\label{eq:adaptive-scale-threshold}
 \frac{M}{6}\le\sigma^2\le\bar\sigma^2,
\end{equation}
then a trial of \Cref{alg:adaptive-utr} either satisfies the stopping
test or satisfies the acceptance test \eqref{eq:acceptance}.
\end{lemma}

\begin{proof}[Proof of \Cref{lem:large-scale}]
If the stopping test holds, the algorithm returns before solving a
subproblem.  Otherwise apply \Cref{lem:common-inexact-tr-step} with
$g=g_k$, $H=H_{k,\sigma}$, $g^+=g_k^+$, and
\[
 a=\sigma\sqrt{h_k}+\frac{\sigma\sqrt\epsilon}{64},\qquad
 R=\frac{\sqrt{h_k}}{4\sigma},\qquad
 \xi=\frac{\sigma\sqrt\epsilon}{64},\qquad
 \delta_0=\delta_+=\delta_g.
\]
The radius and
\eqref{eq:adaptive-scale-threshold} give
$(M/6)\|d\|^3\le\sigma\sqrt{h_k}\|d\|^2/4$.
Thus \eqref{eq:common-function-bound} becomes
\begin{equation}\label{eq:trial-decrease}
    F(x_k)-F(x_k+d)
    \ge\frac14\sigma\sqrt{h_k}\,\|d\|^2
       +\frac12\lambda\|d\|^2-\delta_g\|d\|.
\end{equation}
Similarly, \eqref{eq:common-gradient-bound} yields
\begin{align}
\|g_k^+\|
&\le \sigma\sqrt{h_k}\,\|d\|+\lambda\|d\|
 +2\frac{\sigma\sqrt\epsilon}{64}\|d\|
 +\frac M2\|d\|^2+2\delta_g \notag\\
&\le\left(\frac14+\frac1{128}+\frac3{16}
 +\frac1{128}\right)h_k+\lambda\|d\| \notag\\
&=\frac{29}{64}h_k+\lambda\|d\|
 \le\frac12h_k+\lambda\|d\|. \label{eq:trial-gradient}
\end{align}

If $\|d\|<\sqrt{h_k}/(4\sigma)$, complementarity gives
$\lambda=0$ by \Cref{lem:tr-global-optimality}.  Maximizing
$\delta_g r-\sigma\sqrt{h_k}r^2/4$ over $r\ge0$ and then adding the
two value-oracle errors shows
\[
 \widetilde F_k^+-\widetilde F_k
 \le\frac{\delta_g^2}{\sigma_*\sqrt\epsilon}+\frac\omega4
 \le\omega.
\]
Together with \eqref{eq:trial-gradient}, this proves
\eqref{eq:acceptance-interior}.

If $\|d\|=\sqrt{h_k}/(4\sigma)$, then
$\delta_g\|d\|\le h_k^{3/2}/(1024\sigma)$, so
\eqref{eq:trial-decrease} and the two value-oracle errors give
\[
 \widetilde F_k-\widetilde F_k^+
 \ge\frac{15}{1024}\frac{h_k^{3/2}}\sigma
 +\frac12\lambda\|d\|^2-\frac\omega4.
\]
Condition \eqref{eq:adaptive-boundary-compatibility}, $h_k\ge\epsilon$,
$\sigma\le\bar\sigma$, and $\eta\le1/256$ imply the boundary
acceptance test.  The gradient test follows from
\eqref{eq:trial-gradient}.
Thus every subproblem actually solved with
$\sigma^2\ge M/6$ satisfies the appropriate branch of the acceptance
test, which proves the stated alternative.
\end{proof}

\begin{corollary}[Uniform parameter bound]\label{cor:adaptive-scale-bound}
Suppose \Cref{ass:adaptive-inexact-oracle} holds.  Every trial parameter
generated by \Cref{alg:adaptive-utr} satisfies
$\sigma\le\bar\sigma$, with $\bar\sigma$ defined in
\eqref{eq:adaptive-scale-bar}.  Every inner loop terminates after
finitely many trials.
\end{corollary}

\begin{proof}[Proof of \Cref{cor:adaptive-scale-bound}]
The initial parameter satisfies $\sigma_0\le\bar\sigma$.  Suppose a
trial parameter is at most $\bar\sigma$.  If it is rejected, then
\Cref{lem:large-scale} implies $\sigma<\sigma_*$, so its successor is
smaller than $\gamma_+\sigma_*\le\bar\sigma$.  After an accepted trial
with parameter $\sigma$, the next initial value is
$\max\{\sigma_{\min},\sigma/\gamma_-\}
 \le\max\{\sigma_{\min},\sigma\}$.  Since $\sigma_{\min}\le\sigma_0$,
induction proves the uniform bound.  Finally, repeated rejected-trial
updates reach $\sigma_*$ after finitely many trials, and the first such
trial lies below $\gamma_+\sigma_*\le\bar\sigma$; by
\Cref{lem:large-scale} it cannot be rejected.  Hence every inner loop
is finite.
\end{proof}

With $\inf\varnothing=\infty$, define the accepted-step return index
\begin{equation}\label{eq:adaptive-stopping-time}
 K_\epsilon:=\inf\{k\in\mathbb N_0:
  \text{\Cref{alg:adaptive-utr} returns }x_k\}
 \in\mathbb N_0\cup\{\infty\}.
\end{equation}
For every $k<K_\epsilon$, let $\sigma_k$, $d_k$, and $\lambda_k$
denote the parameter, step, and multiplier of the accepted trial
producing $x_{k+1}$, set $H_k:=H_{k,\sigma_k}$ for its
scale-dependent Hessian, and set $r_k:=\|d_k\|$.
For any finite $N\le K_\epsilon$, define
\begin{equation}\label{eq:adaptive-prefix-sets}
 \mathcal B_N:=\left\{0\le k<N:
  \|d_k\|=\frac{\sqrt{h_k}}{4\sigma_k}\right\},
 \qquad
 \mathcal I_N:=\left\{0\le k<N:
  \|d_k\|<\frac{\sqrt{h_k}}{4\sigma_k}\right\},
\end{equation}
and, for $k\in\mathcal B_N$, set
\begin{equation}\label{eq:z-def}
 z_k:=\frac{\lambda_k\|d_k\|}{h_k}.
\end{equation}

\begin{lemma}[Boundary steps]\label{lem:boundary-count}
Suppose \Cref{ass:adaptive-inexact-oracle} holds.  With the notation in
\eqref{eq:adaptive-stopping-time}--\eqref{eq:z-def}, every finite
$N\le K_\epsilon$ satisfies
\begin{subequations}\label{eq:boundary-bounds}
\begin{align}
 |\mathcal B_N|
 &\le\frac{\bar\sigma(\Delta_F+\omega/4)}
           {\eta\epsilon^{3/2}}
  +\frac{\bar\sigma\omega}{\eta\epsilon^{3/2}}|\mathcal I_N|,
     \label{eq:adaptive-boundary-count}\\
 \sum_{k\in\mathcal B_N}z_k
 &\le\frac{4\bar\sigma(\Delta_F+\omega/4)}
           {\eta\epsilon^{3/2}}
  +\frac{4\bar\sigma\omega}{\eta\epsilon^{3/2}}|\mathcal I_N|.
     \label{eq:adaptive-z-bound}
\end{align}
\end{subequations}
\end{lemma}

\begin{proof}[Proof of \Cref{lem:boundary-count}]
For $k\in\mathcal B_N$, the accepted value test gives a decrease of
the stored proxy of at least
$\eta(h_k^{3/2}/\sigma_k+\lambda_kr_k^2)+\omega$.
For $k\in\mathcal I_N$, that proxy can increase by at most $\omega$.
The two endpoint oracle bounds imply
$\widetilde F_0-\widetilde F_N\le\Delta_F+\omega/4$.
Telescoping the stored proxy therefore gives
\[
 \eta\sum_{k\in\mathcal B_N}
 \left(\frac{h_k^{3/2}}{\sigma_k}+\lambda_kr_k^2\right)
 \le\Delta_F+\frac\omega4+\omega|\mathcal I_N|.
\]
Using $h_k\ge\epsilon$ and $\sigma_k\le\bar\sigma$ proves
\eqref{eq:adaptive-boundary-count}.  For $k\in\mathcal B_N$,
\[
    \lambda_k\|d_k\|^2
    =z_k h_k\|d_k\|
    =\frac{z_k h_k^{3/2}}{4\sigma_k}
    \ge\frac{z_k\epsilon^{3/2}}{4\bar\sigma}.
\]
Summing the multiplier term in
the same telescope proves \eqref{eq:adaptive-z-bound}.
\end{proof}

\begin{lemma}[Gradient-contraction steps]\label{lem:interior-count}
Suppose \Cref{ass:adaptive-inexact-oracle} holds.  With the notation in
\eqref{eq:adaptive-stopping-time}--\eqref{eq:z-def}, every finite
$N\le K_\epsilon$ satisfies
\begin{equation}\label{eq:interior-count}
 |\mathcal I_N|\le |\mathcal B_N|+1+
 \frac{\log_+(h_0/\epsilon)+\sum_{k\in\mathcal B_N}z_k}{\log2}.
\end{equation}
\end{lemma}

\begin{proof}[Proof of \Cref{lem:interior-count}]
For $k\in\mathcal B_N$,
\eqref{eq:acceptance-boundary-gradient} gives
\[
    \frac{h_{k+1}}{h_k}
    \le\max\left\{1,\frac12+z_k\right\},
\]
and therefore
\[
    \left[\log\frac{h_{k+1}}{h_k}\right]_+
    \le z_k.
\]

For $k\in\mathcal I_N$,
\eqref{eq:acceptance-interior} gives the clipped contraction
\[
    h_{k+1}
    =\max\{\|g_{k+1}\|,\epsilon\}
    \le\max\left\{\frac12h_k,\epsilon\right\}.
\]
Thus an index $k\in\mathcal I_N$ with $h_{k+1}>\epsilon$ satisfies
$\log(h_{k+1}/h_k)\le-\log2$.  The complementary case
$h_{k+1}=\epsilon$ means precisely that the clipping floor is active,
or equivalently that $\|g_{k+1}\|\le\epsilon$; the threshold
is $\epsilon$, rather than $\epsilon/2$, because $h_{k+1}$ is clipped
below at $\epsilon$.

Suppose now that $k\in\mathcal I_N$ has $h_{k+1}=\epsilon$ and
$k+1<N$.  Then the algorithm does not return at $x_{k+1}$ and the next
outer iteration accepts a step.  In particular, the stopping test fails
at its accepted parameter $\sigma_{k+1}$.  Since
$\|g_{k+1}\|\le\epsilon$, this failure gives
\[
    \lambda_{\min}\!\left(
       H_{k+1}+\frac{65}{64}\sigma_{k+1}\sqrt\epsilon I
    \right)<0.
\]
Indeed, failure of the Hessian stopping test gives
$\lambda_{\min}(H_{k+1})<-(9/8)\sigma_{k+1}\sqrt\epsilon$, and
$65/64<9/8$.  Since $h_{k+1}=\epsilon$, applying the
positive-semidefiniteness condition in
\eqref{eq:tr-global-optimality} with
$B=H_{k+1}+(65/64)\sigma_{k+1}\sqrt\epsilon I$ forces
$\lambda_{k+1}>0$.  Complementarity then gives
$k+1\in\mathcal B_N$.  Pairing each such floor-hitting index with its
successor is injective, and only $k=N-1$ can remain unpaired.  Therefore
\[
    \bigl|\{k\in\mathcal I_N:h_{k+1}=\epsilon\}\bigr|
    \le|\mathcal B_N|+1.
\]

For the same floor-hitting steps,
$\log(h_{k+1}/h_k)\le0$.  Telescoping $\log h_k$ from $0$ to $N$ and
using $h_N\ge\epsilon$ now gives
\[
    \bigl|\{k\in\mathcal I_N:h_{k+1}>\epsilon\}\bigr|\log2
    \le
    \log\frac{h_0}{\epsilon}
       +\sum_{k\in\mathcal B_N}z_k
    =\log_+\!\left(h_0/\epsilon\right)
       +\sum_{k\in\mathcal B_N}z_k.
\]
Adding the two disjoint interior-step counts proves
\eqref{eq:interior-count}.
\end{proof}

Combining \Cref{lem:boundary-count,lem:interior-count} with
\eqref{eq:adaptive-absorption-load} gives, for every finite prefix of
$N$ accepted nonterminal steps,
\begin{subequations}\label{eq:accepted-complexity}
\begin{align}
 |\mathcal I_N|
 &\le
 \frac{
  (1+4/\log2)
  \bar\sigma(\Delta_F+\omega/4)/(\eta\epsilon^{3/2})
  +1+\log_+(h_0/\epsilon)/\log2
 }{1-\chi},
 \label{eq:adaptive-interior-total}\\
 N
 &\le\frac{\bar\sigma(\Delta_F+\omega/4)}
              {\eta\epsilon^{3/2}}
 +\left(1+\frac{\bar\sigma\omega}{\eta\epsilon^{3/2}}\right)
  |\mathcal I_N|.
 \label{eq:adaptive-accepted-total}
\end{align}
\end{subequations}

\begin{lemma}[Rejected trials]\label{lem:rejected-trials}
Suppose \Cref{ass:adaptive-inexact-oracle} holds.  Let $K_\epsilon$ be the finite
return index in \eqref{eq:adaptive-stopping-time}, and let
$U_\epsilon$ be the number of rejected trials before termination.
Then
\begin{equation}\label{eq:rejected-count}
 U_\epsilon
 \le\frac{\log\gamma_-}{\log\gamma_+}K_\epsilon
 +\frac{\log(\bar\sigma/\sigma_0)}{\log\gamma_+}.
\end{equation}
\end{lemma}

\begin{proof}[Proof of \Cref{lem:rejected-trials}]
For $0\le k\le K_\epsilon$, let $\sigma_k^{\mathrm{init}}$ be the
initial parameter of outer iteration $k$, and let $u_k$ be its number
of rejected trials.  For
$0\le k<K_\epsilon$, the accepted parameter is
$\sigma_k=\sigma_k^{\mathrm{init}}\gamma_+^{u_k}$, so the update rule gives
\[
    \sigma_{k+1}^{\mathrm{init}}
    \ge\frac{\sigma_k^{\mathrm{init}}\gamma_+^{u_k}}{\gamma_-},
\]
and hence
\[
    u_k\log\gamma_+
    \le\log \sigma_{k+1}^{\mathrm{init}}
       -\log \sigma_k^{\mathrm{init}}+\log\gamma_-.
\]
Summing over $0\le k<K_\epsilon$ and using
$\sigma_0^{\mathrm{init}}=\sigma_0$ yields
\begin{equation}\label{eq:accepted-rejection-telescope}
    \sum_{k=0}^{K_\epsilon-1}u_k\log\gamma_+
    \le\log\frac{\sigma_{K_\epsilon}^{\mathrm{init}}}{\sigma_0}
       +K_\epsilon\log\gamma_-.
\end{equation}
In the terminal inner loop, the stopping test returns at parameter
$\sigma_{K_\epsilon}^{\mathrm{init}}
 \gamma_+^{u_{K_\epsilon}}\le\bar\sigma$ by
\Cref{cor:adaptive-scale-bound}.  Thus
\[
    u_{K_\epsilon}\log\gamma_+
    \le\log\frac{\bar\sigma}{\sigma_{K_\epsilon}^{\mathrm{init}}}.
\]
Since
$U_\epsilon=\sum_{k=0}^{K_\epsilon-1}u_k+u_{K_\epsilon}$, adding this
inequality to
\eqref{eq:accepted-rejection-telescope} proves
\eqref{eq:rejected-count}.
\end{proof}

\begin{proof}[Proof of \Cref{thm:total-complexity}]
Combining \Cref{lem:boundary-count,lem:interior-count} and using
\eqref{eq:adaptive-absorption-load} gives
\eqref{eq:accepted-complexity} for every finite nonterminal prefix.
The right-hand side is independent of the prefix length.  Hence
$K_\epsilon=\infty$ would give the same finite upper bound on
$N=|\mathcal B_N|+|\mathcal I_N|$ for arbitrarily large $N$, which is
impossible.  Thus $K_\epsilon<\infty$, and taking $N=K_\epsilon$
establishes the displayed accepted-step bounds.

The final stopping test solves no subproblem.  Therefore
$Q_\epsilon=K_\epsilon+U_\epsilon$, and
\Cref{lem:rejected-trials} gives the trial bound.  At termination, the
current parameter is
at most $\bar\sigma$ by \Cref{cor:adaptive-scale-bound}.  The gradient
oracle error gives the first inequality in
\eqref{eq:return-guarantee}.  Weyl's inequality, the Hessian stopping
test, and \eqref{eq:adaptive-hessian-error} give
\[
 \lambda_{\min}(\nabla^2F(\widehat x))
 \ge-\left(\frac98+\frac1{64}\right)
       \bar\sigma\sqrt\epsilon
 =-\frac{73}{64}\bar\sigma\sqrt\epsilon.
\]
\end{proof}

\subsection{NC--SC minimax results}

\begin{proof}[Proof of \Cref{lem:inner-residual-certificate}]
Strong concavity and Cauchy--Schwarz give the first inequality.
The second follows from $\ell$-smoothness and
$\nabla_yf(x,y^*(x))=0$.
\end{proof}

\begin{proof}[Proof of \Cref{prop:agd-implementation}]
Apply the standard AGD estimate to the smooth strongly convex
objective
\[
 \phi_x(y):=-f(x,y).
\]
See
\citet[Lemma~2]{wang2020improved} or
\citet[Lemma~5]{luo2022finding}.  By
\eqref{eq:inner-residual-certificate}, the initial distance is at most
$R_{\rm in}/\mu$.  It is sufficient to reduce the distance to
$\tau/\ell$, because $\ell$-smoothness then gives the requested
residual.  Combining these facts with
$1-1/\sqrt\kappa\le\exp(-1/\sqrt\kappa)$ proves
\eqref{eq:agd-module-work}.
\end{proof}

\begin{proof}[Proof of \Cref{lem:corrected-errors}]
\smallskip\noindent\emph{Corrected value.}
Fix $x$, write $C=\nabla^2_{yy}f(x,y)$,
$q=\nabla_yf(x,y)$, and $\delta=y-y^*(x)$.  Taylor expansion of the
inner gradient gives
\[
 v:=q-C\delta,
 \qquad
 \|v\|\le\frac\rho2\|\delta\|^2.
\]
If
\[
 \overline C:=\int_0^1\nabla^2_{yy}f(x,y-t\delta)\,dt,
\]
then the fundamental theorem of calculus gives
$q=\overline C\delta$ and hence
$v=(\overline C-C)\delta$.  Both matrices satisfy
$-\ell I\preceq C\preceq-\mu I$ and
$-\ell I\preceq\overline C\preceq-\mu I$, so
\[
 \|v\|\le(\ell-\mu)\|\delta\|.
\]
Taylor expansion of the payoff gives
\[
 P(x)=f(x,y)-q^T\delta+\frac12\delta^TC\delta+R,
 \qquad |R|\le\frac\rho6\|\delta\|^3.
\]
Using $q=C\delta+v$ and completing the square yields
\[
 P(x)-\widehat P(x,y)=R+\frac12v^TC^{-1}v.
\]
Strong concavity gives $\|C^{-1}\|\le1/\mu$.  Moreover,
\[
 \|v\|^2
 \le\frac\rho2\|\delta\|^2
       (\ell-\mu)\|\delta\|,
\]
so
\[
 \frac12\bigl|v^TC^{-1}v\bigr|
 \le\frac{(\kappa-1)\rho}{4}\|\delta\|^3.
\]
Together with the Taylor remainder, this proves
\eqref{eq:corrected-value-error-pf}; its second inequality follows from
$\kappa\ge1$.  If $\|\delta\|\le\mu/\rho$, then the quadratic
Taylor bound on $v$ gives
\[
 \frac12\bigl|v^TC^{-1}v\bigr|
 \le\frac{\rho}{8}\|\delta\|^3.
\]
Combining this estimate with the bound on $R$ proves the
$\kappa$-independent local estimate
\eqref{eq:corrected-value-local-error}.

\smallskip\noindent\emph{Corrected gradient.}
Let
\[
 y_t:=y^*(x)+t\bigl(y-y^*(x)\bigr),
 \qquad t\in[0,1].
\]
Since $f$ is twice continuously differentiable, the fundamental theorem
of calculus applied to the two partial gradients along this segment
gives the exact identities
\begin{align*}
 \nabla_xf(x,y)-\nabla_xf(x,y^*(x))
 &=\int_0^1\nabla^2_{xy}f(x,y_t)
       \bigl(y-y^*(x)\bigr)\,dt,\\
 \nabla_yf(x,y)
 &=\int_0^1\nabla^2_{yy}f(x,y_t)
       \bigl(y-y^*(x)\bigr)\,dt,
\end{align*}
where the second identity uses $\nabla_yf(x,y^*(x))=0$.
Substituting these identities into \eqref{eq:corrected-gradient} and
using \eqref{eq:value-gradient}, we obtain
\begin{align*}
 \widehat g(x,y)-\nabla P(x)
 &=\int_0^1\Bigl[
   \nabla^2_{xy}f(x,y_t)-\nabla^2_{xy}f(x,y)\\
 &\qquad
   +\nabla^2_{xy}f(x,y)
    \bigl(\nabla^2_{yy}f(x,y)\bigr)^{-1}
    \bigl(
      \nabla^2_{yy}f(x,y)-\nabla^2_{yy}f(x,y_t)
    \bigr)
  \Bigr]\\
 &\qquad {}\times\bigl(y-y^*(x)\bigr)\,dt.
\end{align*}
This identity displays the cancellation of the linear inner error.
Hessian Lipschitzness, together with
\[
 \left\|\nabla^2_{xy}f(x,y)
   \bigl(\nabla^2_{yy}f(x,y)\bigr)^{-1}\right\|\le\kappa,
\]
bounds the norm of the integrand by
\[
 (1+\kappa)\rho(1-t)\|y-y^*(x)\|^2.
\]
Integrating over $t\in[0,1]$ proves
\eqref{eq:corrected-gradient-error}.

\smallskip\noindent\emph{Corrected Hessian.}
Retain $C=\nabla^2_{yy}f(x,y)$ and set
\[
 B:=\nabla^2_{xy}f(x,y),\quad
 B_*:=\nabla^2_{xy}f(x,y^*(x)),\quad
 C_*:=\nabla^2_{yy}f(x,y^*(x)).
\]
The inverse identity gives
\[
\begin{aligned}
&\left\|
 \bigl(\nabla^2_{yy}f(x,y)\bigr)^{-1}
 -\bigl(\nabla^2_{yy}f(x,y^*(x))\bigr)^{-1}
 \right\|\\
&\qquad\le\frac{\rho}{\mu^2}\|y-y^*(x)\|.
\end{aligned}
\]
The Schur-complement products satisfy the exact telescoping identity
\begin{align*}
 BC^{-1}B^T-B_*C_*^{-1}B_*^T
 &=(B-B_*)C^{-1}B^T
   +B_*(C^{-1}-C_*^{-1})B^T\\
 &\quad+B_*C_*^{-1}(B-B_*)^T.
\end{align*}
Hessian Lipschitzness, $\|B\|,\|B_*\|\le\ell$, and
$\|C^{-1}\|,\|C_*^{-1}\|\le1/\mu$ bound these three terms by
$\kappa\rho\|y-y^*(x)\|$,
$\kappa^2\rho\|y-y^*(x)\|$, and
$\kappa\rho\|y-y^*(x)\|$, respectively.  Adding the
$xx$-block error gives
\begin{align*}
 \|\widehat H(x,y)-\nabla^2P(x)\|
 &\le \rho\|y-y^*(x)\|
 +2\frac\ell\mu\rho\|y-y^*(x)\|\\
 &\qquad
 +\frac{\ell^2}{\mu^2}\rho\|y-y^*(x)\|\\
 &=(1+\kappa)^2\rho\|y-y^*(x)\|,
\end{align*}
which proves \eqref{eq:corrected-hessian-error}.
\end{proof}

\begin{proof}[Proof of \Cref{cor:ncsc-fixed-oracle}]
The residual--distance certificate and
\eqref{eq:outer-residual-target} give
\[
 \|y-y^*(x)\|
 \le\frac{\sqrt{M_P\epsilon/6}}
          {64(1+\kappa)^2\rho}.
\]
Substitution in \eqref{eq:corrected-hessian-error} gives the Hessian
bound in \eqref{eq:ncsc-fixed-oracle-errors}.  Substitution in
\eqref{eq:corrected-gradient-error} gives
\[
 \|\widehat g-\nabla P\|
 \le\frac{M_P\epsilon}{49152(1+\kappa)^3\rho}
 =\frac{\sqrt2}{12288}
   \left(\frac\kappa{1+\kappa}\right)^3\epsilon
 <\frac\epsilon{8192}.
\]
These bounds are stronger than \eqref{eq:generic-oracle-errors}.
\end{proof}

\begin{proof}[Proof of \Cref{cor:ncsc-adaptive-oracle}]
Equation \eqref{eq:inner-residual-certificate} gives
$\|y-y^*(x)\|\le\delta_{VG}$ or $\delta_H(\sigma)$, respectively.
The first branch of \eqref{eq:ncsc-adaptive-vg-distance} and
\eqref{eq:corrected-gradient-error} give the gradient certificate.
The second branch and \eqref{eq:corrected-value-error-pf} make the
cubic value error at most $\omega/8$.
Equation \eqref{eq:ncsc-adaptive-h-distance} and
\eqref{eq:corrected-hessian-error} give the Hessian certificate.
\end{proof}

\begin{proof}[Verification of
\eqref{eq:ncsc-adaptive-outer-complexity} and
\eqref{eq:ncsc-adaptive-return}]
For \eqref{eq:ncsc-adaptive-explicit-errors}, the interior condition
holds with left-hand side at most $3\omega/4$, while
\[
 \chi=\frac1{16}\left(1+\frac4{\log2}\right)<1
\]
and the boundary condition follows from
$15/1024-1/256=11/1024>5/16384$.
Finally, \Cref{lem:value-geometry} gives the required regularity of
$P$ and \Cref{ass:ncsc} gives $P^*>-\infty$.  Applying
\Cref{thm:total-complexity}, and using
$\|\widehat g_0\|\le\|\nabla P(x_0)\|+\epsilon/256$ to simplify its
initial logarithm, proves
\eqref{eq:ncsc-adaptive-outer-complexity} and
\eqref{eq:ncsc-adaptive-return}.
\end{proof}

\begin{proof}[Proof of \Cref{thm:ncsc-complexity}]
Let $K$ be the number of nonterminal outer iterations and write
$r_k:=\|d_k\|$.
The residual contract and \Cref{cor:ncsc-fixed-oracle} give
\[
 \|g_k-\nabla P(x_k)\|\le\epsilon/256,
 \qquad
 \|H_k-\nabla^2P(x_k)\|
 \le\frac1{64}\sqrt{\frac{M_P}{6}}\sqrt\epsilon.
\]
Thus the outer sequence obeys the fixed-scale repeated-subproblem
procedure in
\Cref{sec:generic-utr}, and \Cref{thm:generic-complexity} applies with
$M=M_P$.  Its logarithmic term uses
$h_0=\max\{\|g_0\|,\epsilon\}$.  Since
\[
 \|g_0\|\le\|\nabla P(x_0)\|+\epsilon/8192,
\]
we have
$\log_+(h_0/\epsilon)
=O(\log_+(\|\nabla P(x_0)\|/\epsilon)+1)$.
Together with
$\sqrt{M_P}=\sqrt{4\sqrt2}\,\kappa^{3/2}\sqrt\rho$, this proves
\eqref{eq:ncsc-return} and \eqref{eq:ncsc-outer-complexity}.

It remains to transfer the work contract of $\mathcal S$.  At the
start of the inner call for $x_{k+1}$, joint $\ell$-smoothness and the
previous residual guarantee give the directly observable warm-start
bound
\[
 \|\nabla_yf(x_{k+1},y_k)\|
 \le\tau_{\rm out}+\ell r_k.
\]
Consequently, apart from the initial call, the logarithmic term in
\eqref{eq:inner-module-work} is bounded by
\[
 1+\log_2\!\left(1+\frac{\ell r_k}{\tau_{\rm out}}\right).
\]
By \eqref{eq:outer-residual-lower},
\[
 \frac{\ell r_k}{\tau_{\rm out}}
 \le64\sqrt3\,r_k\sqrt{\frac{M_P}{\epsilon}}.
\]
Using $\log_2(1+t)=O(1+t^3)$ for $t\ge0$, the logarithm in
each noninitial call is therefore
\[
 O\!\left(
  1+
  \left(r_k\sqrt{\frac{M_P}{\epsilon}}\right)^3
 \right).
\]
By \Cref{lem:cubic-path},
\begin{align*}
 \sum_{k=0}^{K-1}
 \left(r_k\sqrt{\frac{M_P}{\epsilon}}\right)^3
 &=\frac{M_P^{3/2}}{\epsilon^{3/2}}
   \sum_{k=0}^{K-1}r_k^3\\
 &\le\frac{384}{31}\Delta_P\sqrt{M_P}\,\epsilon^{-3/2}
 +\frac{K}{31\sqrt{62}}.
\end{align*}
Summing the module work bounds, adding the initial call, and
substituting the outer bound for $K$ proves
\eqref{eq:ncsc-gradient-complexity}.  Finally,
$\sqrt{M_P}=\sqrt{4\sqrt2}\,\kappa^{3/2}\sqrt\rho$ proves
\eqref{eq:ncsc-gradient-leading}.
\end{proof}

\subsection{Parameter-free residual-module UTR}

\begin{proof}[Source and proof sketch for \Cref{thm:scar-blackbox}]
This is Theorem~5.1 of \citet{lan2026optimal}.  For completeness, its
counting argument is as follows.  An accelerated-restart call at the
current strong-convexity guess $\mu_t$ either detects failure and
replaces the guess by $\mu_t/4$, or reduces the gradient norm by a
factor two.  Once $\mu_t\le\mu$, strong convexity converts the distance
estimate produced by the accelerated subroutine into the latter
gradient reduction.  Thus there are at most
$\lceil\log_4(\mu_0/\mu)\rceil$ failed guesses and at most
$\lceil\log_2(\|\nabla\phi(u_0)\|/\tau)\rceil$ successful reductions.
The accepted guesses remain at least $\mu/4$, so summing the
$O(\sqrt{L/\mu_t})$ cost of the accelerated calls gives
\eqref{eq:scar-original-complexity}.  At the endpoint $L=\mu$, strong
monotonicity, Lipschitz continuity, and Cauchy--Schwarz force
$\nabla\phi(u)-\nabla\phi(v)=L(u-v)$ for every $u,v$.  Taking $v$ to
be the minimizer proves the one-step claim.
\end{proof}

\begin{proof}[Proof of \Cref{cor:scar-inner}]
Strong monotonicity and Lipschitz continuity of $\nabla q_x$ imply
\[
 \mu\|z-y^{\rm in}\|
 \le\|\nabla q_x(z)-\nabla q_x(y^{\rm in})\|
 \le\ell\|z-y^{\rm in}\|.
\]
Consequently, $\mu\le a_x\le\ell$.  The initialization satisfies the
hypotheses of \Cref{thm:scar-blackbox} when $\kappa>1$, and
$\log_4(a_x/\mu)\le\log_4\kappa$.  Substitution in
\eqref{eq:scar-original-complexity} proves the claim in that case.  If
$\kappa=1$, then $a_x=\mu=\ell$ and the endpoint argument following
\Cref{thm:scar-blackbox} gives a numerical constant per call.
\end{proof}

\begin{proof}[Proof of \Cref{lem:pf-asymptotic-oracle}]
Set $\kappa:=\ell/\mu$ and $s:=\sqrt{M_P/6}$.  Choose
$\epsilon_{\rm or}$ to be
\begin{equation}\label{eq:pf-uniform-oracle-onset}
 \min\left\{
  1,\frac{\mu}{\sqrt\rho},
  \left(\frac{\mu^2}{128(1+\kappa)\rho}\right)^{1/3},
  \left(\frac{s\mu}{64(1+\kappa)^2\rho}\right)^{2/3},
  \left(\frac{3\mu^3}{7\rho}\right)^{1/3}
 \right\}.
\end{equation}
The residual--distance inequality in
\eqref{eq:inner-residual-certificate} gives
$\|y-y^*(x)\|\le\epsilon^2/\mu$.  Hence
\eqref{eq:pf-uniform-oracle-onset} ensures
$\epsilon^2/\mu\le\mu/\rho$, and
\Cref{lem:corrected-errors} gives
\begin{align*}
 \|\widehat g-\nabla P\|
 &\le\frac{(1+\kappa)\rho}{2\mu^2}\epsilon^4,\\
 \|\widehat H-\nabla^2P\|
 &\le\frac{(1+\kappa)^2\rho}{\mu}\epsilon^2,\\
 |\widehat P-P|
 &\le\frac{7\rho}{24\mu^3}\epsilon^6.
\end{align*}
The remaining three conditions in
\eqref{eq:pf-uniform-oracle-onset} bound these quantities by
$\epsilon/256$, $s\sqrt\epsilon/64$, and $\epsilon^3/8$,
respectively.  Every coefficient depends only on
$(\ell,\mu,\rho)$, so the choice is uniform over the entire parameter
class and over all $x$ and $y$.
\end{proof}

\begin{proof}[Proof of \Cref{lem:pf-safe-trial}]
Apply \Cref{lem:common-inexact-tr-step} with $F=P$, $M=M_P$,
\[
 a=\sigma\sqrt h+\frac{\sigma\sqrt\epsilon}{64},\qquad
 R=\frac{\sqrt h}{4\sigma},\qquad
 \xi=\frac1{64}\sqrt{\frac{M_P\epsilon}{6}},\qquad
 \delta_0=\delta_+=\frac{(1+\kappa)\rho}{2\mu^2}\epsilon^4.
\]
\Cref{lem:value-geometry} supplies the required Hessian regularity.
At a safe scale, $\xi\le\sigma\sqrt\epsilon/64$, so
$a-\xi\ge\sigma\sqrt h$ and
$a+\xi\le\sigma\sqrt h+2\sigma\sqrt\epsilon/64$.
The radius and \eqref{eq:pf-safe-scale} imply
$(M_P/6)r^3\le\sigma\sqrt h\,r^2/4$.
Thus \eqref{eq:common-function-bound} proves
\eqref{eq:pf-safe-decrease}.

If the radius is active, then
$\sigma\sqrt h\,r^2=h^{3/2}/(16\sigma)$.  By
\eqref{eq:pf-gradient-error}, the last term of
\eqref{eq:pf-safe-decrease} is at most
$h^{3/2}/(1024\sigma)$.  This proves
\eqref{eq:pf-safe-boundary-decrease}.

The common candidate-gradient estimate gives
\[
 \|g^+\|
 \le\sigma\sqrt h\,r+\lambda r
 +2\frac{\sigma\sqrt\epsilon}{64}r
 +\frac{(1+\kappa)\rho}{\mu^2}\epsilon^4
 +\frac{M_P}{2}r^2.
\]
The radius, safe-scale condition, and
\eqref{eq:pf-gradient-error} bound the four terms other than
$\lambda r$ by, respectively,
$h/4$, $h/128$, $h/128$, and $3h/16$.  Their sum is
$116h/256<117h/256$, proving
\eqref{eq:pf-safe-gradient}.
\end{proof}

\begin{proof}[Proof of \Cref{lem:pf-accepted-accounting}]
Denote the current corrected gradient, accepted step, trial scale, and
multiplier of accepted step $k$ by $g_k,d_k,\sigma_k$, and
$\lambda_k$, and set
$h_k:=\max\{\|g_k\|,\epsilon\}$ and $r_k:=\|d_k\|$.  For the stated
prefix, define
\[
 \mathcal B_N:=\left\{0\le k<N:
   r_k=\frac{\sqrt{h_k}}{4\sigma_k}\right\},
 \qquad
 \mathcal I_N:=\left\{0\le k<N:
   r_k<\frac{\sqrt{h_k}}{4\sigma_k}\right\}.
\]
Every $k\in\mathcal B_N$ gives a true decrease of at least
\[
 \eta\left(\frac{h_k^{3/2}}{\sigma_k}
       +\lambda_kr_k^2\right),
\]
because the two corrected-value errors sum to at most
$\epsilon^3/4$.  Every $k\in\mathcal I_N$ can increase the true
value by at most $5\epsilon^3/4$.  Telescoping and using
$h_k\ge\epsilon$, $\sigma_k\le\bar\sigma$ give
\begin{subequations}\label{eq:pf-boundary-budgets}
\begin{align}
 |\mathcal B_N|
 &\le\frac{\bar\sigma\Delta_P}
          {\eta\epsilon^{3/2}}
 +\frac{5\bar\sigma\epsilon^{3/2}}{4\eta}|\mathcal I_N|,
       \label{eq:pf-boundary-count}\\
 \sum_{k\in\mathcal B_N}\frac{\lambda_kr_k}{h_k}
 &\le\frac{4\bar\sigma\Delta_P}
          {\eta\epsilon^{3/2}}
 +\frac{5\bar\sigma\epsilon^{3/2}}{\eta}|\mathcal I_N|.
       \label{eq:pf-multiplier-budget}
\end{align}
\end{subequations}
Here we used, for $k\in\mathcal B_N$,
$\lambda_kr_k^2=(\lambda_kr_k/h_k)h_k^{3/2}/(4\sigma_k)$.

The accepted gradient tests give
\[
 h_{k+1}\le
 \begin{cases}
  \max\{h_k/2,\epsilon\},&k\in\mathcal I_N,\\
  \max\{h_k/2+\lambda_kr_k,\epsilon\},&k\in\mathcal B_N.
 \end{cases}
\]
Thus every $k\in\mathcal I_N$ not hitting the floor $\epsilon$
decreases $\log(h_k/\epsilon)$ by at least $\log2$, whereas
$k\in\mathcal B_N$ can
increase it by at most $\lambda_kr_k/h_k$.

It remains to count floor-hitting indices in $\mathcal I_N$.  If such
an index satisfies $k+1<N$, then at the next accepted outer iteration
$h=\epsilon$ and the stopping test fails.  The gradient branch of that
test already holds; otherwise $h>\epsilon$, contrary to the floor
event.  Hence the Hessian branch must fail.  To show that the next
accepted solution is on the boundary, let
\[
 A=H+\frac{\sigma\sqrt\epsilon}{64}I
      +\sigma\sqrt\epsilon I.
\]
Failure of the stopping test gives
$-\lambda_{\min}(A)>7\sigma\sqrt\epsilon/64$.
Applying \eqref{eq:tr-global-optimality} with $B=A$ therefore forces
$\lambda>0$, and complementarity forces
$r=\sqrt\epsilon/(4\sigma)$.
Rejected trials only increase $\sigma$ and retest stopping, so each
nonterminal floor event is injectively followed by an index in
$\mathcal B_N$; only the last prefix index can remain unpaired.
Therefore
\begin{equation}\label{eq:pf-interior-potential}
 |\mathcal I_N|
 \le |\mathcal B_N|+1+
 \frac{
  \log_+(\|g_0\|/\epsilon)
  +\sum_{k\in\mathcal B_N}\lambda_kr_k/h_k
 }{\log2}.
\end{equation}
Substituting \eqref{eq:pf-boundary-budgets} into
\eqref{eq:pf-interior-potential} and applying
\eqref{eq:pf-absorption-condition} bounds
$N=|\mathcal B_N|+|\mathcal I_N|$ by the right-hand side of
\eqref{eq:pf-accepted-bound}, uniformly in $N$.
\end{proof}

\begin{proof}[Proof of \Cref{thm:pf-scar-utr}]
Set
\[
 s:=\sqrt{M_P/6},
 \qquad
 \bar\sigma:=2s.
\]
Choose the class-uniform threshold no larger than
$\epsilon_{\rm or}$ in \Cref{lem:pf-asymptotic-oracle} and so that
\begin{equation}\label{eq:pf-small-epsilon-conditions}
\begin{gathered}
 \epsilon\le s,\qquad
 \frac{5\bar\sigma\epsilon^3}
 {4\eta\epsilon^{3/2}}
 \left(1+\frac4{\log2}\right)\le\frac12,\\
 \frac1{s\sqrt\epsilon}
 \left(\frac{(1+\kappa)\rho}{2\mu^2}\epsilon^4\right)^2
 \le\frac34\epsilon^3,\\
 \left(\frac1{128}-\eta\right)
 \frac{\epsilon^{3/2}}{\bar\sigma}
 \ge\frac54\epsilon^3.
\end{gathered}
\end{equation}
For example, take
\begin{equation}\label{eq:pf-class-uniform-onset}
 \epsilon_0:=\frac12\min\left\{
  \epsilon_{\rm or},s,
  \left(
   \frac{\eta}{5s(1+4/\log2)}
  \right)^{2/3},
  \left(
   \frac{3s\mu^4}{(1+\kappa)^2\rho^2}
  \right)^{2/9}
 \right\}.
\end{equation}
For every $0<\epsilon\le\epsilon_0$, the term containing
$\eta/[5s(1+4/\log2)]$ implies both the absorption and boundary-acceptance
conditions in \eqref{eq:pf-small-epsilon-conditions}; the last term
implies its interior-increase condition.  Thus this single choice,
which depends only on $(\ell,\mu,\rho)$, ensures all the displayed
requirements.

Starting from
$\sigma=\epsilon$, a rejected trial doubles $\sigma$.  Unless the
algorithm has already returned at its stopping or zero-step test, the
first nonzero trial with $\sigma\ge s$ has $\sigma\le\bar\sigma$.
By \Cref{lem:pf-safe-trial}, its corrected-gradient test passes.  An
interior trial can have a small true increase because of the remaining
gradient error.  Minimizing the right-hand side of
\eqref{eq:pf-safe-decrease} over $r\ge0$ shows that this increase is at
most
\[
 \frac1{\sigma\sqrt h}
 \left(\frac{(1+\kappa)\rho}{2\mu^2}\epsilon^4\right)^2
 =O(\epsilon^{15/2}).
\]
Here we used $h\ge\epsilon$ and $\sigma\ge s$.  This is
at most $3\epsilon^3/4$ by
\eqref{eq:pf-small-epsilon-conditions}; after
adding the two value errors, the interior value test passes.  On a
boundary trial,
\eqref{eq:pf-safe-boundary-decrease} and the two value errors imply
\[
 \widehat P-\widehat P^+
 \ge\frac{h^{3/2}}{128\sigma}
     +\frac14\lambda r^2-\frac{\epsilon^3}{4}.
\]
Because $h\ge\epsilon$ and $\sigma\le\bar\sigma$, the right-hand side
is at least $\eta(h^{3/2}/\sigma+\lambda r^2)+\epsilon^3$ by
\eqref{eq:pf-small-epsilon-conditions}.  Hence a safe-scale trial
either stops or is
accepted, and no generated scale exceeds $\bar\sigma$.

The absorption condition in \Cref{lem:pf-accepted-accounting} follows
from \eqref{eq:pf-small-epsilon-conditions}.  Consider an arbitrary
finite trial prefix, and let $K_N$ and $U_N$ be its numbers of accepted
and rejected nonterminal trials.  Applying
\Cref{lem:pf-accepted-accounting} to its accepted-step subsequence
bounds $K_N$ by the right-hand side of
\eqref{eq:pf-accepted-bound}, independently of the prefix length.
Starting from $\sigma=\epsilon$, each rejection doubles $\sigma$ and
each acceptance decreases it by at most a factor of two.  Since every
generated scale is at most $\bar\sigma$, this gives
\[
 U_N\le K_N+\log_2(\bar\sigma/\epsilon)+1.
\]
Thus the length $K_N+U_N$ of every finite trial prefix has one common
finite upper bound.  An infinite run would have arbitrarily long
finite prefixes, a contradiction.  The algorithm therefore
terminates.  Applying the same bounds to the complete finite run, with
$K$ accepted and $U$ rejected trials, proves the same estimate with
$\log_+(\|g_0\|/\epsilon)$ in place of the first logarithm in
\eqref{eq:pf-outer-complexity}, because $Q\le K+U+1$; the extra unit
accounts for a terminal solved subproblem with $d=0$.  At the initial
tuple, \eqref{eq:pf-gradient-error} gives
$\|g_0\|\le\|\nabla P(x_0)\|+\epsilon/256$.  Hence replacing
$\log_+(\|g_0\|/\epsilon)$ by
$\log_+(\|\nabla P(x_0)\|/\epsilon)$ costs at most a universal
additive constant, which proves \eqref{eq:pf-outer-complexity}.

If the explicit stopping test returns, then
\eqref{eq:pf-gradient-error}, $\sigma\le2s$, and
\eqref{eq:pf-hessian-error} give
\[
 \|\nabla P(x)\|\le\frac{257}{256}\epsilon,
 \qquad
 \lambda_{\min}(\nabla^2P(x))
 \ge-\frac{145}{64}s\sqrt\epsilon
 \ge-\sqrt{M_P\epsilon},
\]
where $145<64\sqrt6$.  If the robust
subproblem returns $d=0$, global optimality of the differentiable
quadratic model gives $g=0$, $h=\epsilon$, and
\[
 H+\frac{\sigma\sqrt\epsilon}{64}I
   +\sigma\sqrt\epsilon I
 \succeq0.
\]
Together with \eqref{eq:pf-oracle-bounds} and $\sigma\le2s$, this gives
\[
 \|\nabla P(x)\|\le\frac\epsilon{256},
 \qquad
 \lambda_{\min}(\nabla^2P(x))
 \ge-\frac{131}{64}s\sqrt\epsilon
 \ge-\sqrt{M_P\epsilon},
\]
where $131<64\sqrt6$.  A constant rescaling of the internal
tolerance converts the first inequality in \eqref{eq:pf-return} to
$\|\nabla P(\widehat x)\|\le\epsilon$.

Finally,
$\sqrt{M_P}=\sqrt{4\sqrt2}\,\kappa^{3/2}\sqrt\rho$ gives the stated
leading outer dependence.
\end{proof}

\begin{proof}[Proof of \Cref{cor:pf-scar-inner-complexity}]
Every accepted boundary step decreases the true value, while every
accepted interior step can increase it by at most $5\epsilon^3/4$.
There are at most $Q$ accepted steps, so every accepted outer point
satisfies
\[
 P(x)-P^*\le\Delta_P+\frac54Q\epsilon^3.
\]
Since $\nabla P$ is $(1+\kappa)\ell$-Lipschitz, the standard descent
inequality and \eqref{eq:pf-gradient-error} give $h\le G_Q$.

At a candidate point, the SCAR module is warm-started from an inner point with
residual at most $\epsilon^2$.  The joint $\ell$-smoothness and the
trust-region radius give the initial candidate residual bound
\[
 \epsilon^2+\ell\|d\|
 \le\epsilon^2+\frac{\ell\sqrt{G_Q}}{4\epsilon}.
\]
Apply \Cref{cor:scar-inner} with target $\epsilon^2$ to every candidate
call and separately to the initial call at $(x_0,y_{-1})$.  There is
at most one candidate call per solved subproblem, so summation gives
\eqref{eq:pf-inner-complexity}.
\end{proof}

\subsection{Adaptive fixed-radius SCAR trust region}

\begin{proof}[Proof of \Cref{lem:scar-persistent-half}]
The AR estimate used in the proof of
\citet[Theorem~5.1]{lan2026optimal}, which follows from their
Theorem~4.1 and remains valid when $L=\mu$, states that a call with the
current guess $\nu$ and an admissible line-search estimate uses at
most
\[
 4+C_1\sqrt{10L/\nu}
\]
gradient evaluations and returns a point $v$ satisfying
\[
 \|\nabla\phi_j(v)\|
 \le\frac\nu2\|u-u_j^*\|,
 \qquad u_j^*:=\arg\min_u\phi_j(u).
\]
If $\nu\le\mu$, strong convexity gives
$\|u-u_j^*\|\le\|\nabla\phi_j(u)\|/\mu$, so the returned residual is
at most $R/2$.  Therefore a failed halving test implies $\nu>\mu$.
Every failure divides $\nu$ by four, and once $\nu\le\mu$ no later
function in the sequence can cause another failure.  This proves
\eqref{eq:persistent-scar-failures}.  Moreover, throughout the process
$\nu\ge\mu/4$.  Hence each attempted stage costs at most
\[
 4+C_1\sqrt{40\kappa}
 \le(4+2\sqrt{10}\,C_1)\sqrt\kappa.
\]
There are $N$ successful attempts and at most the number of failed
attempts in \eqref{eq:persistent-scar-failures}, which proves
\eqref{eq:persistent-scar-count}.  The line-search statement is the
uniform estimate in \Cref{thm:scar-blackbox}.  Since all $\phi_j$ have
the same global smoothness bound $L$, a retained estimate remains an
admissible initialization after the objective changes.
\end{proof}

\begin{proof}[Proof of \Cref{lem:working-residual-invariant}]
Lipschitz continuity of the full gradient gives
\[
 \|\nabla_yf(x+d,y)\|
 \le\|\nabla_yf(x,y)\|+\ell\|d\|
 \le(1+c_0)\ell R(\sigma).
\]
After $N_{\rm w}$ successful halvings this is at most
$c_0\ell R(\sigma)$ by \eqref{eq:working-half-count}.  At the same
$x$, one halving reduces the old residual bound to
$c_0\ell R(\sigma)/2=c_0\ell R(2\sigma)$.
\end{proof}

\begin{proof}[Proof of \Cref{lem:fixed-radius-oracle-accuracy}]
Because
$c_0\kappa\sqrt\epsilon\,\mathcal A_\epsilon\to0$ as
$\epsilon\downarrow0$, there is a threshold depending only on
$(\ell,\mu,\rho)$ below
which \eqref{eq:working-distance-bound} and
$\sigma\ge\mathcal A_\epsilon^{-1}$ imply
$\|y-y^*(x)\|\le\mu/\rho$.  Substituting
\eqref{eq:working-distance-bound} into
\eqref{eq:corrected-value-local-error} then gives
\eqref{eq:fixed-radius-proxy-error}, proving part~(i).

For part (ii), \Cref{lem:corrected-errors} and
\eqref{eq:working-distance-bound} imply
\begin{align*}
 \|\widehat g-\nabla P\|
 &\le\frac{(1+\kappa)\rho}{2}
       \frac{c_0^2\kappa^2\epsilon}{\sigma^2},\\
 \|\widehat H-\nabla^2P\|
 &\le(1+\kappa)^2\rho
       \frac{c_0\kappa\sqrt\epsilon}{\sigma}.
\end{align*}
Using $\kappa\ge1$, $M_P=4\sqrt2\kappa^3\rho$, and
$M_P\le6\sigma^2$, the right-hand sides are at most
\[
 \frac{3c_0^2}{2\sqrt2}\epsilon,
 \qquad
 3\sqrt2c_0\sigma\sqrt\epsilon,
\]
respectively.  With $c_0=2^{-12}$ these are bounded by the first two
quantities in \eqref{eq:working-safe-errors}.  For the value bound, the
ratio of \eqref{eq:fixed-radius-proxy-error} to
$\epsilon^{3/2}/\sigma$ is at most $7c_0^3/(1024\sqrt2)$, which is
smaller than $1/4096$.  This proves the third inequality.

For part (iii), $\mu\le a_x\le\ell$ and
\eqref{eq:inner-residual-certificate} give
\[
 \|y_{\rm v}-y^*(x)\|
 \le c_0\kappa\epsilon^{3/2}.
\]
The corrected-gradient and corrected-Hessian errors are, uniformly in
$x$, at most
\[
 \frac{(1+\kappa)\rho c_0^2\kappa^2}{2}\epsilon^3,
 \qquad
 (1+\kappa)^2\rho c_0\kappa\epsilon^{3/2},
\]
respectively.  The first is at most $\epsilon/1024$ for sufficiently
small $\epsilon$.  For the second, note that
$\epsilon\mathcal A_\epsilon\to0$ and
$\sigma\ge\mathcal A_\epsilon^{-1}$, so
the displayed Hessian error is at most
$\sigma\sqrt\epsilon/32$ for all sufficiently small $\epsilon$.
This proves \eqref{eq:validator-errors}.  All estimates used only the
common constants $(\ell,\mu,\rho)$, so the thresholds are uniform over
the stated class.
\end{proof}

\begin{proof}[Proof of \Cref{lem:adaptive-fixed-radius-safe}]
Take $\epsilon_3\le\min\{\epsilon_1,\epsilon_2\}$ from
\Cref{lem:fixed-radius-oracle-accuracy}.
Let
\[
 m(d):=\widehat g^Td+\frac12d^T
 \left(\widehat H+\frac{65}{64}\sigma\sqrt\epsilon I\right)d.
\]
Apply \Cref{lem:common-inexact-tr-step} with $F=P$, $M=M_P$,
\[
 a=\frac{65}{64}\sigma\sqrt\epsilon,\qquad
 R=R(\sigma),\qquad
 \xi=\frac{\sigma\sqrt\epsilon}{64},\qquad
 \delta_0=\frac{\epsilon}{1024},\qquad
 \delta_+=0,
\]
and $g^+=\nabla P(x+d)$.  The model part of the common estimate gives
\begin{equation}\label{eq:fixed-radius-model-bound}
 m(d)\le-\frac\lambda2\|d\|^2.
\end{equation}

At a safe scale, \eqref{eq:working-safe-errors}, the Taylor bound
\eqref{eq:function-taylor}, and
\eqref{eq:fixed-radius-model-bound} yield, with $r=\|d\|$,
\begin{equation}\label{eq:fixed-radius-taylor-model}
\begin{aligned}
 P(x+d)-P(x)
 \le{}&-\frac\lambda2r^2
 -\frac{65}{128}\sigma\sqrt\epsilon\,r^2
 +\frac{\epsilon}{1024}r
 +\frac{\sigma\sqrt\epsilon}{128}r^2
 +\frac{M_P}{6}r^3.
\end{aligned}
\end{equation}
For a boundary step, $r=\sqrt\epsilon/(4\sigma)$ and
$M_P r/6\le\sigma\sqrt\epsilon/4$.  Rearranging
\eqref{eq:fixed-radius-taylor-model} gives
\[
 P(x)-P(x+d)
 \ge\frac14\sigma\sqrt\epsilon\,r^2
     +\frac\lambda2r^2-\frac{\epsilon}{1024}r,
\]
which is exactly \eqref{eq:safe-boundary-true-decrease}.  The two
endpoint value errors in \eqref{eq:working-safe-errors} then give
\begin{equation}\label{eq:safe-boundary-proxy-decrease}
 \widehat P(x,y)-\widehat P(x+d,y^+)
 \ge\frac{61}{4096}\frac{\epsilon^{3/2}}{\sigma}
     +\frac\lambda2R(\sigma)^2
 >\frac{\epsilon^{3/2}}{256\sigma}.
\end{equation}
Thus the boundary step is accepted.

Now suppose $r<R(\sigma)$.  Complementarity gives $\lambda=0$.
The common candidate-gradient estimate gives
\begin{align*}
 \|\nabla P(x+d)\|
 &\le\frac{\epsilon}{1024}
 +\frac{\sigma\sqrt\epsilon}{64}r
 +\frac{65}{64}\sigma\sqrt\epsilon\,r
 +\frac{M_P}{2}r^2\\
 &\le\frac{1+4+260+192}{1024}\epsilon
 =\frac{457}{1024}\epsilon.
\end{align*}
Applying the positive-semidefiniteness part of
\eqref{eq:tr-global-optimality} with
$B=\widehat H+(65/64)\sigma\sqrt\epsilon I$ and using Hessian
Lipschitz continuity likewise imply
\begin{align*}
 \lambda_{\min}(\nabla^2P(x+d))
 &\ge-\frac{65}{64}\sigma\sqrt\epsilon
      -\frac{1}{64}\sigma\sqrt\epsilon-M_P r\\
 &\ge-\frac{162}{64}\sigma\sqrt\epsilon
 =-\frac{81}{32}\sigma\sqrt\epsilon.
\end{align*}
The validation errors in \eqref{eq:validator-errors} therefore give
\[
 \|\widehat g_{\rm v}\|
 \le458\epsilon/1024<\epsilon/2,
 \qquad
 \lambda_{\min}(\widehat H_{\rm v})
 \ge-82\sigma\sqrt\epsilon/32
 >-(21/8)\sigma\sqrt\epsilon.
\]
This proves (ii).  Finally, if the validation test succeeds, the same
error bounds give
\[
 \|\nabla P(x^+)\|
 \le\epsilon/2+\epsilon/1024<\epsilon,
\]
and
\[
 \lambda_{\min}(\nabla^2P(x^+))
 \ge-\left(\frac{21}{8}+\frac1{32}\right)
       \sigma\sqrt\epsilon
 =-\frac{85}{32}\sigma\sqrt\epsilon.
\]
This proves (iii).
\end{proof}

\begin{proof}[Proof of \Cref{thm:adaptive-fixed-radius-scar}]
Set
\begin{equation}\label{eq:fixed-radius-scale-cap}
 \sigma_*:=\sqrt{M_P/6},
 \qquad
 \overline\sigma:=2\sigma_*=2\sqrt{M_P/6}.
\end{equation}
These quantities are positive because \Cref{ass:ncsc} assumes
$\rho>0$.
Choose $\epsilon_0\in(0,1)$ small enough that
\Cref{lem:fixed-radius-oracle-accuracy,lem:adaptive-fixed-radius-safe}
apply and
\begin{equation}\label{eq:initial-scale-below-safe}
 \sigma_0=\mathcal A_\epsilon^{-1}<\sigma_*.
\end{equation}
The last condition is possible because
$\mathcal A_\epsilon\to\infty$ as $\epsilon\downarrow0$ and the
common constants give $M_P>0$.  Every restriction imposed here and in
the two cited lemmas depends only on $(\ell,\mu,\rho)$; hence the same
$\epsilon_0$ works for every payoff in the class and every
initialization.

The scale is doubled only after a rejected boundary trial or a failed
interior validation.  By \Cref{lem:adaptive-fixed-radius-safe}, once a
generated scale is at least $\sigma_*$, a boundary trial is accepted
and an interior trial terminates the method.  The first generated scale
at least $\sigma_*$ is smaller than $2\sigma_*$, so no scale exceeds
$\overline\sigma$.  Counting the doublings needed to cross from
$\mathcal A_\epsilon^{-1}$ to $\sigma_*$ proves
\begin{equation}\label{eq:fixed-radius-doubling-count}
 J\le
 \left\lceil\log_2(\sigma_*\mathcal A_\epsilon)\right\rceil.
\end{equation}

It remains to bound accepted boundary steps, including those at unsafe
scales.  The corrected value is not a one-sided bound on $P$, so we
must account explicitly for changes in the stored proxy.  At scale
$\sigma$, \eqref{eq:working-proxy-bound} gives
\[
 |\widehat P-P(x)|\le\mathcal E_\epsilon(\sigma).
\]
Every accepted boundary step decreases the stored proxy by at least
$\epsilon^{3/2}/(256\sigma)$.  The candidate tuple becomes the next
stored tuple without recomputation.  A failure leaves $x$ unchanged,
doubles the scale, and recomputes the current tuple.  At such a refresh
from $\sigma$ to $2\sigma$, the upward jump of the stored proxy is at
most
\begin{equation}\label{eq:proxy-refresh-jump}
 \mathcal E_\epsilon(\sigma)
 +\mathcal E_\epsilon(2\sigma).
\end{equation}
The failure scales form a subset of
$\{2^j\sigma_0:j\ge0\}$.  Extending the sum to the whole set gives
\begin{equation}\label{eq:cubic-refresh-series}
 \sum_{j\ge0}\left(2^{-3j}+2^{-3(j+1)}\right)=\frac97.
\end{equation}
The initial proxy is at most
$P(x_0)+\mathcal E_\epsilon(\sigma_0)$, and every final stored proxy is
at least $P^*-\mathcal E_\epsilon(\sigma_0)$ because the scale never
decreases.  Therefore, for every finite prefix with $N$ accepted
boundary steps, enumerate those steps chronologically and let
$\sigma_k$ denote the scale of the $k$th accepted boundary trial.  Then
\begin{align}
 \frac{\epsilon^{3/2}}{256}
 \sum_{k=0}^{N-1}\frac1{\sigma_k}
 \le{}&\Delta_P
 +\frac{23}{1536}c_0^3\rho\kappa^3
      \epsilon^{3/2}\mathcal A_\epsilon^3.
 \label{eq:fixed-radius-proxy-telescope}
\end{align}
Indeed, before substituting \eqref{eq:fixed-radius-proxy-error}, the
proxy-error multiplier $23/7$ is the sum of two endpoint copies and the
refresh sum in \eqref{eq:cubic-refresh-series}.
Since every $\sigma_k\le\overline\sigma$,
\eqref{eq:fixed-radius-proxy-telescope} implies
\begin{align}
 N\le{}&256\overline\sigma\Delta_P\epsilon^{-3/2}
 +\frac{23}{6}\overline\sigma
   c_0^3\rho\kappa^3\mathcal A_\epsilon^3.
 \label{eq:fixed-radius-boundary-count}
\end{align}

This uniform finite-prefix bound also proves termination.  Suppose, to
the contrary, that the loop is infinite.  At any fixed unsafe scale,
infinitely many accepted boundary steps would contradict
\eqref{eq:fixed-radius-boundary-count}.  Hence a nonaccepted trial must
eventually occur.  Under the nontermination hypothesis, such a trial
either fails the interior validation or rejects a boundary candidate,
and therefore doubles the scale.  After finitely many such doublings a
safe scale is generated.  At a safe scale, an interior trial terminates,
a boundary rejection is impossible, and an infinite sequence of
accepted boundary steps again contradicts
\eqref{eq:fixed-radius-boundary-count}.  This contradiction proves
termination.
Every nonterminal trial is either an accepted boundary or a failure,
and the final trial is interior, so $Q=K+J+1$.  Combining
\eqref{eq:fixed-radius-doubling-count} and
\eqref{eq:fixed-radius-boundary-count}, now with $N=K$, gives the more
explicit estimate
\begin{align}
 Q\le{}&256\overline\sigma\Delta_P\epsilon^{-3/2}
 +\frac{23}{6}\overline\sigma
   c_0^3\rho\kappa^3\mathcal A_\epsilon^3
 \notag\\
 &+\left\lceil\log_2(\sigma_*\mathcal A_\epsilon)\right\rceil+1.
 \label{eq:fixed-radius-explicit-outer-count}
\end{align}
Substituting $\overline\sigma=2\sqrt{M_P/6}$ proves
\eqref{eq:fixed-radius-leading-outer}; the coefficient of its leading
term is $512/\sqrt6$.

Finally, \eqref{eq:validator-true-return} and
$\sigma\le\overline\sigma$ give
\eqref{eq:fixed-radius-return}.  Since
\[
 \left(\frac{85}{16\sqrt6}\right)^2=\frac{7225}{1536},
\]
replacing the internal tolerance by $(1536/7225)\epsilon$ gives the
standard curvature threshold $-\sqrt{M_P\epsilon}$; its gradient
bound is only strengthened.

It remains to count inner-gradient evaluations.  Define, only for this
accounting argument,
\begin{align}
 N_{\rm init}
 &:={}
 \left\lceil
  \log_2\max\left\{1,
   \frac{\|\nabla_yf(x_0,y_{-1})\|}
        {c_0a_0R(\sigma_0)}
  \right\}
 \right\rceil,
 \label{eq:initial-half-count}\\
 N_{\rm val}
 &:={}
 \left\lceil
  \log_2\max\left\{1,
   \frac{\kappa\mathcal A_\epsilon}{4\epsilon}
  \right\}
 \right\rceil.
 \label{eq:validation-half-count}
\end{align}
The initial sequence of halving calls needs at most
$N_{\rm init}$ successful calls to reach
$c_0a_0R(\sigma_0)$.  Every trust-region trial uses exactly
$N_{\rm w}$ successful calls to construct its working candidate.
Every failure uses one additional successful call to refresh the
current tuple at the doubled scale.  Hence these operations require at
most
\[
 N_{\rm init}+N_{\rm w}Q+J
\]
successful halvings.

It remains to count validation.  There are at most $J$ failed
validators and one terminal validator.  Immediately before validation,
the candidate satisfies the working invariant, so its residual is at
most
\[
 c_0\ell R(\sigma)
 =\frac{c_0\ell\sqrt\epsilon}{4\sigma}.
\]
The target in \eqref{eq:validator-residual} is at least
$c_0\mu\epsilon^{3/2}$.  Since
$\sigma\ge\sigma_0=\mathcal A_\epsilon^{-1}$, the ratio between the
starting residual and the validation target is at most
$\kappa\mathcal A_\epsilon/(4\epsilon)$.  Therefore each validator
requires at most $N_{\rm val}$ additional successful halvings.  The
total number of successful halving calls is consequently bounded by
\[
 N_{\rm w}Q+J+(J+1)N_{\rm val}+N_{\rm init}.
\]

Apply \Cref{lem:scar-persistent-half} to the entire sequence of inner
objectives $q_x=-f(x,\cdot)$.  They share the same $(\mu,\ell)$, and
the initial secant obeys $\mu\le a_0\le\ell$.  Thus all failed
curvature guesses together contribute at most
$\lceil\log_4\kappa\rceil$ attempted stages, not one such term per
outer trial.  Computing the initial secant and at most $J+1$
validation secants costs no more than $2(J+2)$ additional
inner-gradient evaluations.  Consequently,
\begin{align}
 N_y\le{}&(4+2\sqrt{10}\,C_1)\sqrt\kappa
 \Bigl[
  N_{\rm w}Q+J+(J+1)N_{\rm val}+N_{\rm init}
  +\lceil\log_4\kappa\rceil
 \Bigr]
 +2(J+2).
 \label{eq:fixed-radius-inner-explicit}
\end{align}

Finally, $N_{\rm w}$ is a universal integer, and
\[
 \begin{aligned}
 J&=O_{\rm fixed}(\log\mathcal A_\epsilon+1),\\
 N_{\rm val}&=O_{\rm fixed}(\log(1/\epsilon)),\\
 N_{\rm init}&=O_{\rm fixed}(\log(1/\epsilon)).
 \end{aligned}
\]
The last relation uses $a_0\ge\mu$ and
$R(\sigma_0)=\sqrt\epsilon\,\mathcal A_\epsilon/4$.  Substituting
\eqref{eq:fixed-radius-explicit-outer-count} into
\eqref{eq:fixed-radius-inner-explicit} proves
\eqref{eq:fixed-radius-inner-leading}.  Finally,
$\sqrt{\kappa M_P}=\Theta(\kappa^2\sqrt\rho)$ gives the equivalent
leading dependence stated in the theorem.
\end{proof}

\section{Scope of the parameter-free guarantees}
\label{app:parameter-free-scope}
This section distinguishes the class-uniform asymptotic guarantees
from a residual-only certificate that would be uniform without a
positive common strong-concavity lower bound.  It also clarifies the
additional bookkeeping used by the fixed-radius construction.

\subsection{The residual-only obstruction}
\label{app:scar-limitation}

A parameter-free inner method such as SCAR can return a point satisfying
the prescribed and directly verifiable condition
$\|\nabla_yf(x,y)\|\le\tau$.  This fact alone does not provide a
certificate for the value or gradient of
$P(x)=\max_y f(x,y)$ that is uniform over problem classes having no
positive common lower bound on the strong-concavity modulus.  The issue
is not merely a loss in the complexity bound: an outer stopping test
based only on a residual-corrected oracle can be false.

The following scalar family shows that this difficulty persists even
for the corrected value \eqref{eq:corrected-value-pf}.  Let
\[
 c:=\frac12,
 \qquad
 \psi(t):=\frac{t^2}{2}+c(1-\cos t),
 \qquad
 A:=\frac{1-c}{2\pi^2},
\]
and, for $0<\mu\le1/4$, define
\begin{equation}\label{eq:scar-counterexample}
 f_\mu(x,y):=A\tanh(x/A)
       \bigl(1-\cos(2\pi\sqrt\mu\,y)\bigr)
       -\psi(2\sqrt\mu\,y-1).
\end{equation}
Because $1-c\le\psi''(t)\le1+c$,
\[
 \partial_{yy}^2 f_\mu(x,y)
 =4A\pi^2\mu\tanh(x/A)\cos(2\pi\sqrt\mu\,y)
   -4\mu\psi''(2\sqrt\mu\,y-1)
 \le-2(1-c)\mu=-\mu.
\]
Thus $f_\mu(x,\cdot)$ is globally $\mu$-strongly concave.  Its
second derivatives are uniformly bounded for $0<\mu\le1/4$.  Its third
derivatives are also uniformly bounded: the four distinct block types
have bounds $12/A^2$, $2\pi/A$, $\pi^2$, and $A\pi^3+c$,
respectively.  Hence the joint gradient-smoothness and
Hessian-Lipschitz constants admit finite upper bounds independent of
$\mu$.  Moreover,
\[
 P_\mu(x)\ge f_\mu\left(x,\frac1{2\sqrt\mu}\right)
 =2A\tanh(x/A)\ge-2A,
\]
so the value function is bounded below.
At $x=0$ and $y=1/(2\sqrt\mu)$,
$-\partial_{yy}^2f_\mu(x,y)=4\mu\psi''(0)=6\mu$.
Consequently, the largest valid global strong-concavity modulus lies
between $\mu$ and $6\mu$ and genuinely tends to zero along the family.

At $x=0$, the exact maximizer is
$y_\mu^*(0)=1/(2\sqrt\mu)$.  At the point $y=0$, however,
\[
 |\partial_yf_\mu(0,0)|
 =2\sqrt\mu\,|\psi'(-1)|
 =(2+\sin 1)\sqrt\mu\longrightarrow0,
 \qquad
 \widehat g_\mu(0,0)=0,
 \qquad
 \widehat H_\mu(0,0)=0.
\]
In contrast, Danskin's theorem gives
\[
 P_\mu'(0)
 =\partial_xf_\mu\left(0,\frac1{2\sqrt\mu}\right)
 =1-\cos\pi=2.
\]
The corrected value is wrong by a fixed amount as well.  Namely,
$P_\mu(0)=0$, whereas
\begin{equation}\label{eq:scar-counterexample-value}
 \widehat P_\mu(0,0)
 =-\psi(-1)+\frac{\psi'(-1)^2}{2\psi''(-1)}
 \approx0.0647,
\end{equation}
independently of $\mu$.  Therefore a residual tending to zero does not
uniformly control either the corrected value or the corrected gradient
over the union of these parameter classes.  The smoothness constants
remain bounded while the best strong-concavity modulus tends to zero.

Consequently, choosing SCAR as a residual module removes the inner
solver's dependence on $\ell$ and $\mu$, but the residual alone does
not yield a single certificate valid across all positive
strong-concavity moduli.  Nor does a prescribed finite sequence of
positive tolerances repair the issue.  For any finite list
$\tau_0,\ldots,\tau_J>0$, choose
\[
 0<\mu<\min\left\{\frac14,
  \left(\frac{\min_j\tau_j}{2+\sin 1}\right)^2\right\}
\]
and repeat each inner call at the outer point $x=0$, initialized at
$y=0$.  Every residual test is then already satisfied and may return
the same false outer derivatives.
This observation covers any predetermined finite number of
tolerance-halving phases.  It does not claim that indefinite
refinement cannot eventually move the inner iterate; it shows that a
finite prescribed refinement depth is not a certificate.

There is no contradiction with the class-uniform result in
\Cref{thm:pf-scar-utr}.  Once a positive common value of $\mu$ is
fixed, the condition
\[
 \epsilon^2<(2+\sin 1)\sqrt\mu
\]
excludes $y=0$ as a legal return of an inner call with residual
tolerance $\epsilon^2$.  More generally, every legal return satisfies
\[
 \|y-y_\mu^*(x)\|\le\frac{\epsilon^2}{\mu}.
\]
Thus one threshold may be chosen uniformly after the common constants
$(\ell,\mu,\rho)$ have been fixed.  What fails is a threshold uniform
over their union as the common lower bound $\mu$ tends to zero.

SCAR also does not return a certified lower bound on $\mu$, a certified
inner optimality gap, or an interval containing $P(x)$.  Its returned
smoothness quantity is a line-search estimate satisfying the tested
local inequalities, not a certified global upper bound on $\ell$.
Likewise, a successful residual reduction does not certify that the
current strong-convexity guess is below the true $\mu$.  A fully
observable finite-accuracy adaptive minimax certificate therefore
needs additional
information, such as a certified lower bound on $\mu$, a certified
inner-gap bound, genuine intervals containing $P$, $\nabla P$, and
$\nabla^2P$, compactness plus a dual-gap certificate, or additional
problem structure.  \Cref{thm:pf-scar-utr} instead gives an unknown
class-uniform threshold $\epsilon_0(\ell,\mu,\rho)$; it makes neither
a claim uniform as $\mu\downarrow0$ nor an a posteriori claim that the
algorithm can recognize entry into this regime.

\subsection{Additional qualifications for the fixed-radius method}
\label{app:fixed-radius-scope}

The class-uniform quantifier in
\Cref{thm:adaptive-fixed-radius-scar} has the same nonobservability
caveat: the algorithm uses no unknown problem parameter, but it cannot
observe whether $\epsilon\le\epsilon_0(\ell,\mu,\rho)$.  The proof uses
the positive common lower bound on $\mu$ to convert working and
validation residuals into distances and corrected-oracle errors.

The counterexample in \Cref{app:scar-limitation} concerns prescribed
absolute residual tolerances.  It does not, by itself, rule out a
threshold uniform over $\mu\downarrow0$ for the present fixed-radius
construction, whose working and validation targets are scaled by
observable secants.  Accordingly, \Cref{thm:adaptive-fixed-radius-scar}
asserts uniformity only after $(\ell,\mu,\rho)$ have been fixed and
makes no claim either way about a stronger union-uniform theorem.

The fixed-radius construction handles the function-value difficulty
in a specific, nonstandard way.  It never asserts that
$\widehat P$ is a lower or an upper bound on $P$.  Unsafe accepted
boundary steps need not decrease the true value function.  Instead,
the proof telescopes the stored corrected proxy and explicitly pays
for every proxy jump caused by a scale refresh; see
\eqref{eq:proxy-refresh-jump}--\eqref{eq:fixed-radius-proxy-telescope}.
This argument relies on three implementation details: the scale is
nondecreasing, an accepted candidate inner point and corrected value
are retained without further inner refinement, and a rejected candidate
is discarded before the current tuple is refreshed.  Without these
bookkeeping rules, the displayed telescoping argument no longer
applies without additional error accounting.

The high-accuracy validator is invoked only after an interior
fixed-radius solution.  Its auxiliary residual
$c_0a_x\epsilon^{3/2}$ controls inner accuracy only; it is not used as
the trust-region scale or as the Hessian stopping threshold.  Under
the theorem's class-uniform small-$\epsilon$ hypotheses, the validator
rules out false termination, but it does not reveal whether those
hypotheses hold.  A doubling trick in $\epsilon$ does not by itself
identify the unknown onset of this asymptotic regime.  The present
analysis therefore does not furnish an observable all-$\epsilon$
certificate.  Establishing one would require a new certificate or
argument, for example a certified lower bound on $\mu$, a certified
inner-gap or value-function interval, compactness with a computable dual
gap, or other problem-specific structure.

\bibliographystyle{unsrtnat}
\bibliography{ref}

\end{document}
